%%%%%%%%%%%%%%%%%%%%%%%%%%%%%%%%%%%%%%%%%%%%%%%%%%%%%%%%%%%%%%%%%%%%%%%%%%%%%%%
%% Autonomous Robots (Springer) submission.
%% Class: Springer Nature sn-jnl, formatting option [iicol] as the journal requires.
%% Reference style: sn-apa -- the journal specifies author-year citations, an alphabetised
%% reference list, italic journal titles and DOIs given as full links.
%%
%% Authors, affiliations, ORCIDs (comment block below the author list, for the submission
%% system), funding, competing interests, CRediT contributions, both availability statements with
%% the archive landing page, and the appendices are included. The frozen review package and its
%% submission checklist distinguish current evidence from earlier public archive versions.
%% The standard scope for every quantity is
%% the configuration census at five seed sets where five exist; see the remark "Scope of the
%% quantities in this paper" in Sect. 5.5, and Appendix C for the per-cell run counts, which differ
%% by task.
%%
%% Build: pdflatex main; bibtex main; pdflatex main; pdflatex main
%%%%%%%%%%%%%%%%%%%%%%%%%%%%%%%%%%%%%%%%%%%%%%%%%%%%%%%%%%%%%%%%%%%%%%%%%%%%%%%
\documentclass[pdflatex,iicol,sn-apa]{sn-jnl}

\usepackage{graphicx}
\usepackage{multirow}
\usepackage{amsmath,amssymb,amsfonts}
\usepackage{amsthm}
\usepackage[utf8]{inputenc}
\usepackage{textcomp}
\usepackage{booktabs}
\usepackage{url}

\theoremstyle{remark}
\newtheorem{remark}{Remark}

\raggedbottom

\begin{document}

\title[What Determines a Simulation Benchmark Ranking?]{What Determines a Simulation Benchmark
Ranking? Structural Blindness, Finite Configuration Grids, and Evaluation Budget in
Simulation-Based Policy Comparison}

%% Author order as supplied, with Yi Sui last as corresponding author. Converted from the
%% elsarticle-style block supplied (\ead/\fnref/\cortext/\address) into the Springer sn-jnl macros;
%% \author* marks the corresponding author and \affil* their affiliation.
\author[1]{\fnm{Jun} \sur{Ji}}\email{junji@qdu.edu.cn}
\author[2]{\fnm{Bowen} \sur{Tan}}\email{btanab@connect.ust.hk}
\author[1]{\fnm{Yi} \sur{Li}}\email{ly2005@qdu.edu.cn}
\author[1]{\fnm{Xiaolei} \sur{Zhang}}\email{zhangxiaolei@qdu.edu.cn}
\author[3]{\fnm{Shengjie} \sur{Guo}}\email{guosj@emails.imau.edu.cn}
\author*[1]{\fnm{Yi} \sur{Sui}}\email{suiyi@qdu.edu.cn}

\affil*[1]{\orgdiv{College of Computer Science and Technology}, \orgname{Qingdao University},
\orgaddress{\city{Qingdao}, \postcode{266071}, \country{China}}}

\affil[2]{\orgname{The Hong Kong University of Science and Technology},
\orgaddress{\city{Hong Kong SAR}, \country{China}}}

\affil[3]{\orgname{Inner Mongolia Agricultural University},
\orgaddress{\city{Hohhot}, \country{China}}}

%% ORCIDs, for entry into the submission system. The sn-jnl class does define \orcid{<url>}, but it
%% renders an Orcidlayout logo from Orcidlogo.eps, which the template package does not ship, so the
%% ids are recorded here rather than typeset. They are also in declarations_for_interface.md.
%%   Jun Ji            0000-0003-3194-2183
%%   Bowen Tan         0009-0007-0554-9261
%%   Yi Li             0000-0002-4185-3152
%%   Xiaolei Zhang     0000-0002-0122-4554
%%   Shengjie Guo      0009-0004-6852-4836
%%   Yi Sui            0009-0001-8081-5183

\abstract{Simulation benchmarks guide robot foundation policy selection, but calibration and repeated
evaluation address different sources of ranking uncertainty. We audit SIMPLER across two simulator
implementations, two robot embodiments and eight policy configurations.
First, replaying 98 public demonstrations reveals weakly observed controller directions. On a
64-configuration census with five seed sets, varying these directions changes one estimated policy
gap by 0.120 and changes whether individual settings support a declaration. Separate re-calibration
on same-source demonstrations favours a different operating point, without establishing an error
in the original fit. Second, the benchmark evaluates finite populations of 24 to 300 configurations,
and its implementations can enumerate different populations. A complete census removes configuration
sampling error while leaving policy randomness and differences between simulator settings.
Third, the reference 72-episode budget resolves only one or two of four reported orderings,
depending on the variance model. Across 17 policy pairs, Holm-adjusted analyses retain nine
single-setting declarations and five declarations over sampled parameter sets. The four additional
abstentions are not confirmed corrections of real-robot errors.
These findings connect calibration, configuration coverage and evaluation budget to policy
comparisons. The accompanying audit protocol makes the evaluation population, operating point,
calibration loss, search domain and policy randomness explicit. It supports simulation-sensitivity
diagnosis without providing a guarantee of real-world ordering.}

\keywords{Policy evaluation, Sim-to-real transfer, Identifiability, Partial identification,
Benchmark reproducibility, Robot manipulation}

\maketitle

%%=============================================================================%%
\section{Introduction}\label{sec1}

A simulator calibrated against real demonstrations is increasingly treated as evidence about which
robot policy is better. Simulation suites that mirror real evaluation scenes
\citep{li2024simpler} are used to compare foundation policies at a scale no physical laboratory can
match, and their published tables are read as rankings: this policy beats that one on this task.
The appeal is obvious. A physical evaluation of a single policy on a single task costs hours of
robot time and is difficult to reproduce; a simulated one costs minutes and is, in principle,
exactly repeatable.

The standard question asked of such a suite is how well it correlates with reality, and it is
usually answered with a correlation coefficient or a rank agreement over a handful of policies.
This paper asks a question that comes before that one, and that can be answered without a robot:

\begin{quote}
Given a simulation benchmark's stated calibration protocol, the configurations it enumerates and
the evaluation budget it uses, which ranking claims are supported, and which depend on choices
that those measurements leave unresolved?
\end{quote}

The distinction matters because a correlation study treats the simulator's output as a measurement
with some error attached, whereas our question asks whether the output is a well-defined function
of the inputs at all. If two simulator configurations are equally consistent with every piece of
calibration data and they disagree about which policy wins, then no amount of correlation analysis
repairs the ranking: the evidence does not contain the answer. This possibility must be distinguished
from insufficient \emph{statistical power}, which additional independent evaluations can address.
Our audit measures both sources of ambiguity and separates their effects rather than assuming
that every unstable declaration is evidence of non-identifiability.

We answer the question on a widely used sim-to-real evaluation suite for manipulation, across two
embodiments (a WidowX bridge setup and a Google-robot fractal setup), eight policy configurations
across three model families --- Octo \citep{octo2024}, OpenVLA \citep{kim2024openvla} and RT-1
\citep{brohan2023rt1} --- and a complete enumeration of every task's configuration grid. Three
findings, each independent of the others.

\paragraph{The calibration protocol is structurally blind (Sect.~\ref{sec4}).}
The protocol identifies some dynamics parameters and is, over the range it excites, uninformative
about others.
Replaying 98 real demonstrations, we show analytically and then empirically that a common scaling of
the joint PD gains and the value of the torque limit leave the replay trajectory unchanged to within
magnitudes we quantify below --- the
first because the free-space response depends on the ratio $d/k$ up to inertial terms that these
slow, teleoperated demonstrations do not excite, the second because commanded
torque never saturates along demonstrated velocities. Sect.~\ref{subsec4_2} states both as
insensitivity over a measured domain rather than as exact invariance, and quantifies the residual. The measured effect is 0.141~mm on one
simulator stack and 0.207~mm on a second, independent one for the gain scaling. For the torque limit
the replayed trajectory is bitwise identical on 97 of the 98 episodes, on both stacks; the single
exception is the same demonstration on each, differing by 0.149~$\mu$m on one stack and 0.050~mm on
the other. That pattern is what the saturation argument predicts --- invariance holds exactly where
the commanded torque stays inside the limit, and one demonstration evidently approaches it --- so we
state the invariance as near-exact rather than exact. These are structural invariances of the
protocol, not flat regions of a noisy objective, so no improvement in fitting technique removes
them.

\paragraph{The benchmark is a finite configuration population, not a sample (Sect.~\ref{sec5}).}
The episode index determines the initial state modulo the size of the task's configuration grid ---
24 configurations for three bridge tasks, 64 for one port's eggplant task, 300 for the fractal
coke-can task. Episodes beyond that count are exact repeats, so the common practice of running more
episodes stops buying information about configurations at a threshold most users never learn. Two
widely used ports of the same suite enumerate \emph{different} grids, so a number computed on one is
not comparable to a number computed on the other. Enumerating the grid drives configuration sampling
error to exactly zero, and what remains is instructive: policy-seed noise, implementation-build
drift, and the suite's own averaging over four robot \emph{texture} variants, across which one
policy's success rate ranges over 0.400 --- that is, 40 percentage points, the unit in which every
effect size in this paper is quoted. These three are not commensurable and we do not rank them: they
are a standard deviation, a paired shift and a range over four settings. The texture parameter
changes the colour of the robot model, carries no physics whatsoever, and is invisible to replay
calibration for exactly that reason; the averaging over it is part of the documented protocol, and
what we contribute is its magnitude and the observation that it makes the published number a rate
marginal over the four variants, which is a different estimand from the conditional rate a
single-variant reproduction measures.

\paragraph{The published protocol does not resolve the rankings it reports (Sect.~\ref{sec6}).}
We follow the reference protocol, including its pseudo-random-number lifecycle, and reproduce the
published success rates to a mean absolute difference of 0.038 across eight cells, and on the second
embodiment to $+0.003$ and $-0.040$. The differences are one-signed --- six of the eight cells fall
below the published value, two match it, none exceeds it, mean $-0.038$ --- and we report that
because our own first reproduction did not show it: that run re-seeded the policy every episode,
which inflated its variance and let a noisier estimator look like a closer match
(Sect.~\ref{subsec6_2}). We do not convert those eight cells into a platform constant; they share
tasks, policies and seeds. What licenses the rest of the paper is narrower and sufficient: every
effect we report is a difference between two policies measured in the same configuration on one
pipeline and one build, where an offset common to both policies cancels. That last condition is not
decoration: Sect.~\ref{subsec5_5} measures a change of build shifting $\hat\Delta$ itself by
$-0.109$, so a build change is \emph{not} a common offset, and we never compare across one. At the
protocol's own 72-episode budget, one of the four policy orderings it reports is statistically
resolvable under the variance model the protocol's own randomisation implies, and two under a
coarser model that treats each run rather than each configuration as the replication unit; we report
both and argue for the former (Sect.~\ref{subsec6_3}). On carrot the published difference is $+0.014$
and ours is $+0.028$ with interval $[-0.039, +0.094]$ under either: a number whose sign its own
budget does not determine, which is precisely the failure mode this paper is about.

Having established that ranking claims are underdetermined from three independent directions, we ask
what can be done about it (Sects.~\ref{sec7} and~\ref{sec8}). The natural remedy is a set-valued
verdict: declare an ordering only when it holds for every simulator parameter the calibration data
cannot exclude. We measure that criterion's reach, and report both halves of the answer. It is
conservative as designed on synthetic ground truth, where the point criterion's error rate
\emph{rises} with episode count to 0.88, and it stabilises verdicts on real replicate data. At the
reference protocol's three seed sets it disagrees with point calibration on 3 of 17 policy pairs, and
at a single seed set on 2 of the 13 pairs one seed set can evaluate at all --- but when we spend more
seeds, two of the three disagreements turn out to have been our own insufficient power, leaving
one. And on a pair we selected from the 5 of 62 published pairs whose real and simulated means point
opposite ways, the criterion over the invisible fibre \emph{declares the ordering the published real
means reverse}. Whether that is a model gap or noise in a real evaluation published without
intervals we cannot say, and Sect.~\ref{subsec8_4} argues the two readings support the same
conclusion: a declaration from this criterion carries no warranty about the robot. Its reach stops
at the parameter family the calibration data can constrain.

\subsection{Contributions}\label{subsec1_1}

This paper contributes three empirical findings about robot-policy evaluation. Each connects a
specific property of the evaluation protocol to a measurable consequence for a policy comparison.

\begin{description}
\item[C1.] \textbf{Calibration sensitivity and policy-comparison sensitivity can diverge.}
Replay of 98 demonstrations is weakly informative about common controller-gain scaling and an
inactive torque limit over the tested ranges (Sect.~\ref{sec4}). On the ManiSkill3 eggplant census,
one policy pair's estimated gap spans 0.120 across settings, with configuration coverage, seed
budget and inference build held fixed (Table~\ref{tab6a}). The table identifies the separate replay
and policy data sources. A torque change shifts the cross-family comparisons by 0.104--0.146
(mean 0.121)
(Sect.~\ref{subsec7_3}); this is a separate measurement. These are changes in gaps and declarations,
not demonstrated reversals of sign-determined rankings.
\item[C2.] \textbf{The configuration population determines what a benchmark success rate estimates.}
The evaluated tasks contain 24 to 300 index-determined configurations. Their enumeration eliminates
configuration sampling error for that population, while policy randomness remains. Different ports
can enumerate different populations, and averaging over four robot texture variants estimates a
different quantity from evaluating one variant; one policy's rate spans 0.400 across those variants
(Sect.~\ref{sec5}). This separates uncertainty reduced by repeated runs from differences requiring
an explicit choice of evaluation target.
\item[C3.] \textbf{Ranking declarations depend on budget, analysis model and operating point.}
The reference 72-episode budget resolves one or two of four reported orderings, depending on the
variance model (Sect.~\ref{sec6}). At the released per-pair budgets, Holm adjustment across 17 bridge
pairs retains nine single-setting declarations and five sampled-set declarations, giving four
additional abstentions under the stated inference assumptions (Sect.~\ref{subsec7_2}). Separately,
re-calibration on same-source demonstrations favours a different controller setting. Moving there
changes one point verdict while both completed set envelopes abstain; the paired change in the
policy gap itself is unresolved (Sect.~\ref{subsec7_7}). This is a scope-bounded re-calibration result,
not evidence that the original fit was erroneous or that four real-world errors were corrected.
\end{description}

\noindent The set-valued rule is a classical sensitivity tool, not a new statistical test. Its
component-interval assumptions, the construction of its parameter set and its lack of a real-world
ordering guarantee are stated separately in Sects.~\ref{sec3} and~\ref{sec8}. The practical output is
an audit protocol that identifies the evaluation population, calibration operating point and loss,
search domain, inference build and policy-randomness budget before interpreting a ranking.

\noindent The appendices carry the reproducibility assets: the original reference stack running
headless on a GPU-less host through a purpose-built Vulkan compatibility layer, the
configuration-census protocol, and append-only provenance logs for every episode reported here.

\subsection{Audit trail}\label{subsec1_2}

Appendix~\ref{secE} records four earlier claims withdrawn after additional checks, with the
measurements and protocol corrections that changed them. The main text reports the final claims
and retains the relevant sensitivity analyses; the appendix preserves the development history for
readers assessing selection and reproducibility.

\subsection{Scope}\label{subsec1_3}

We make no claim about which policy is actually better on any real robot, and we ran no hardware
ourselves. Where real-robot ground truth is needed --- in Sect.~\ref{subsec8_4}'s delimiting
experiment --- we use the success rates published alongside the benchmark for the same policies and
task, which are the reference values the benchmark itself is validated against. Two public
evaluation services that would have provided independent hardware results were unavailable during
this work: one endpoint was offline throughout and the other's queue allocation could not be
confirmed. Our claims are about what the simulated evidence determines, which is a question about
the benchmark rather than about the robot.

%%=============================================================================%%
\section{Related work}\label{sec2}

\paragraph{Simulation-based evaluation of manipulation policies.}
Recent evaluation suites construct simulated scenes that mirror real evaluation setups, calibrate
the simulator against real demonstrations, and report success rates intended to be comparable with
real-robot numbers \citep{li2024simpler}. Parallel efforts pursue automated real-robot evaluation
\citep{zhou2025autoeval} and distributed evaluation across laboratories
\citep{atreya2025roboarena}. All are motivated by the cost and irreproducibility of physical
evaluation. What validates each of them differs and the difference matters here: SIMPLER reports
rank correlation and a mean maximum rank violation against published real rates, AutoEval automates
the real evaluation itself and so validates against its own hardware runs, and RoboArena aggregates
distributed real evaluations rather than validating a simulator at all. Our work takes such a suite as its
object of study rather than as a tool, and asks not how closely it agrees with reality but which of
its own outputs are determined by its own inputs.

We should be precise about the suite we study, because it anticipates more than a correlation
analysis. \citet{li2024simpler} discuss the weaknesses of Pearson correlation for this purpose,
note that small real differences make rank agreement unstable, and propose a mean maximum rank
violation to measure ordering agreement directly. What we add is not the observation that
correlation is a poor summary. It is a per-pair analysis conditioned on a stated calibration
protocol and a stated evaluation budget: which individual ranking claims that evidence determines,
where the parameter directions the protocol cannot see enter, and how the answer moves with the
operating point and the observation unit.

\paragraph{Confidence statements for simulated and real evaluation.}
The closest neighbours are methods that attach a finite-sample statement about \emph{real}
performance to a simulated or partly-simulated evaluation. \citet{suresim2025} use a small set of
paired real and simulated evaluations with prediction-powered inference to give intervals for real
performance, and \citet{scape2026} give scenario-conditioned predictions with conformal
calibration. \citet{polaris2025} belongs beside these but is not the same kind of thing: it is a
real-to-sim reconstruction and evaluation pipeline validated by empirical correlation with real
evaluations, not a method that carries a finite-sample interval, and we separate it here so the
grouping does not imply a guarantee it does not claim. Our estimand is the simulation benchmark's
pairwise policy difference, conditional on an explicit configuration population, operating point
and evaluation protocol. This addresses a complementary question: which simulation-side choices
need to be fixed or reported before a real-performance correction can be interpreted? Paired real
and simulated evaluations are needed to compare real-performance interval methods empirically;
we do not have those data. The present comparison concerns simulation-side declaration sensitivity,
and neither its parameter envelope nor its calibration set is a substitute for real-data calibration.

\paragraph{Identifiability and system identification.}
That not every parameter of a dynamical model can be recovered from a given excitation is classical,
and structural identifiability analysis is standard in system identification \citep{ljung1999system,
bellman1970structural}. Robot calibration practice generally acknowledges it in the form of
excitation design: trajectories are chosen to make the parameters of interest observable
\citep{swevers1997optimal}. The gap we address is that demonstration-replay calibration does not get
to choose its excitation --- the demonstrations are whatever humans happened to teleoperate, they
are free-space and slow, and the resulting blindness is inherited rather than designed. We treat the
blind directions not as a nuisance to be reduced but as a set to be propagated, which connects to
the next literature.

\paragraph{Partial identification and set-valued inference.}
When data do not pin down a parameter, one can report the set of values consistent with the data and
the corresponding set of implied conclusions, rather than a point estimate and a confidence interval
around it \citep{manski2003partial, tamer2010partial}. Model-confidence-set procedures provide the
specific device we use: invert a test to obtain the collection of models the data cannot
distinguish, then require a conclusion to hold across that collection \citep{hansen2011model}. Our
contribution is not methodological novelty here but instantiation and measurement --- we construct
the compatible set for a real calibration protocol, propagate it to policy rankings, and quantify
changes in declarations and the limits of extrapolation to real performance.

\paragraph{Randomising what the simulator gets wrong.}
A large literature responds to simulator inaccuracy by randomising over it during training, so that
the learned policy is robust to whichever value is right --- over visual appearance
\citep{tobin2017domain} and over dynamics \citep{peng2018sim}. Our subject is the mirror image of
that programme. Randomisation is a technique for making a \emph{policy} insensitive to parameters
the modeller cannot pin down; we ask what remains of an \emph{evaluation} when those same
parameters cannot be pinned down. The two interact in a way worth stating plainly, because one of the
terms that survives a configuration census on our data is a purely visual one: it is not news that
visual variation matters for transfer, which is precisely the premise of domain randomisation. Nor
is the aggregation over it undocumented --- the suite reports it as part of its protocol
\citep{li2024simpler}. What we add is the magnitude, and the observation that averaging over such a
variant makes the published number a rate \emph{marginal over} the variants rather than the rate at
any one of them, so a single-variant reproduction targets a different quantity
(Sect.~\ref{subsec5_5}). That is a reporting and estimand point, not an argument against
randomisation.

\paragraph{Reproducibility and statistical practice in learned-policy evaluation.}
Sensitivity of reinforcement learning results to random seeds, and the resulting unreliability of
small-sample comparisons, is documented and has produced concrete recommendations on reporting
\citep{henderson2018deep, agarwal2021deep, colas2018how}. Our Sect.~\ref{sec5} contributes two items
this literature does not cover for simulated robot benchmarks: the initial-state grid is finite and
enumerable, so the configuration axis is a population rather than a sample and its sampling error
can be driven to exactly zero; and the \emph{implementation build} --- the port branch plus the
inference stack version --- behaves as a parameter of the same order as the ones being
studied. The first changes which bootstrap is correct; the second means that even a perfectly
specified bootstrap understates the uncertainty of a number copied from a published table.

%%=============================================================================%%
\section{Problem setup}\label{sec3}

\subsection{Notation}\label{subsec3_1}

Let $z$ denote the simulator's physical parameters (joint PD gains, torque limits, contact friction
and density), and let $\pi$ denote the \emph{implementation build}: the port branch, the renderer
path and the policy inference stack version. A task provides a finite set $\mathcal{C}$ of initial
configurations. A policy $i$ run from configuration $c$ under parameters $z$ and build $\pi$ with
internal random source $\sigma$ either succeeds or fails, written
$s_i(c, z, \sigma; \pi) \in \{0,1\}$. The simulated difference between two policies is
\begin{equation}\label{eq:delta}
\begin{split}
\Delta^{\mathrm{sim}}_{ij}(z; \pi) = \frac{1}{|\mathcal{C}|} \sum_{c \in \mathcal{C}}
\big[ \, &\mathbb{E}_\sigma\, s_i(c, z, \sigma; \pi) \\[-2pt]
         &- \mathbb{E}_\sigma\, s_j(c, z, \sigma; \pi) \, \big],
\end{split}
\end{equation}
and $\Delta^{\mathrm{real}}_{ij}$ denotes the corresponding real-robot difference. Three properties
of this definition carry the weight of Sect.~\ref{sec5} and are stated here because they are easy to
overlook: the sum over $\mathcal{C}$ is over a \emph{population}, not a sample; $\sigma$ is
degenerate for deterministically decoding policies, so for those the benchmark value is measurable
exactly; and $\mathcal{C}$ itself depends on $\pi$, so two builds with different grids yield two
different estimands that cannot be paired per episode.

\subsection{The compatible rule, and why we do not call it a confidence set}\label{subsec3_2}

Calibration observes real demonstrations and scores a candidate $z$ by a replay loss $L(z)$. Rather
than reporting the minimiser $\hat z$, we invert a one-sided test at level $\alpha$: a parameter
setting is \emph{compatible} if its per-demonstration error is not significantly worse than the
best,
\begin{equation}\label{eq:cset}
C_\alpha = \Big\{ z : \mathrm{lo}_{1-\alpha}\big( \mathbb{E}[L(z) - L(\hat z)] \big) \le 0 \Big\},
\end{equation}
where $\mathrm{lo}_{1-\alpha}(\cdot)$ is the paired-bootstrap $(1-\alpha)$ lower confidence bound.
The threshold is exactly zero: the definition admits no engineering-equivalence margin, and if one
is wanted it has to be introduced explicitly and justified.

\paragraph{This is an empirical rule, not a $1-\alpha$ confidence set, and we were wrong to call it
one.} Each test is valid \emph{for a fixed reference}. Ours is not fixed: $\hat z$ is the argmin of
the same demonstrations the test then uses, so it is the most favourable candidate in this draw and
candidates near it are rejected more often than $\alpha$ allows. Inverting a test does not handle
that selection. The size of the effect is easy to see under a null of $K$ candidates with identical
risk, where every candidate is a true minimiser: at $K = 50$ and $M = 98$, matching our grid and
demonstration count, a nominated true minimiser is retained about $53\%$ of the time against a
nominal $95\%$ (\texttt{scripts/check\_selection\_bias.py}). Increasing the bootstrap count does not
help, because the problem is the reference rather than the resampling. We therefore make \textbf{no
coverage claim} for Eq.~\eqref{eq:cset} as computed, and read it as what it is: a reproducible
empirical rule for which settings this loss and this grid cannot separate from the best one.

\paragraph{A variant that removes the selection dependence, and what it says.} Selecting the reference on one half
of the demonstrations and testing on the other makes selection and inference independent. On the
same null that variant retains a true minimiser $92.4\%$ of the time --- most of the gap closed,
with a few points of residual shortfall from the percentile bootstrap itself at $M/2 \approx 49$, so
we call it a large improvement rather than an exact procedure. We report it alongside the
same-sample rule wherever the distinction could matter, and Sect.~\ref{subsec4_4} shows that the
paper's substantive conclusion --- that the benchmark's nominal setting is excluded --- holds under
it on both simulator stacks and across five random splits. Directions in which the loss is invariant
(Sect.~\ref{subsec4_2}) enter $C_\alpha$ at every magnitude regardless of where the minimum lies,
which is a stronger statement than ``the data are consistent with a range of values'': along such an
axis there is no range, the axis is unconstrained. Parameters the replay cannot observe at all ---
contact friction and object density never appear in a free-space trajectory --- are likewise
unconstrained, and we flag them separately because their compatibility is a property of the
protocol's coverage rather than of its precision.

Two objects follow from this and must not be conflated, because in our data they differ.
$C_\alpha$ is the set Eq.~\eqref{eq:cset} defines, computed against the loss minimiser over a stated
search domain; Sect.~\ref{subsec4_4} computes it and finds that it \emph{excludes} the simulator's
nominal setting. $F_{\mathrm{nom}}$, defined in Eq.~\eqref{eq:cset2}, is the fibre of settings that
leave the replay loss unchanged \emph{relative to nominal} --- the invisible axes through the point
at which the benchmark actually runs. Our verdicts range over $F_{\mathrm{nom}}$, because every
published rate is computed at nominal, and we say so wherever it matters.

\subsection{The decision problem}\label{subsec3_3}

Given a parameter set $S$, an evaluation budget and a policy pair, we must either declare an ordering
or abstain. Two criteria are compared throughout:
\begin{itemize}
\item \textbf{point calibration}, standard practice: evaluate at a single setting and declare
$i \succ j$ when the interval for $\hat\Delta_{ij}$ there excludes zero;
\item \textbf{the union bound}: declare only when every $z \in S$ agrees in sign, using the
interval $\big[\min_{z} \mathrm{lo}(z),\ \max_{z} \mathrm{hi}(z)\big]$.
\end{itemize}

\noindent The second needs no multiplicity correction across the settings in $S$, and the reason is
worth stating because it is the only place this paper gets something for free. Throughout,
$[\mathrm{lo}(z), \mathrm{hi}(z)]$ is an \emph{equal-tailed two-sided} $95\%$ interval, so each
one-sided tail carries $\alpha_1 = 0.025$; we write $\alpha_1$ for that tail and not for the nominal
$0.05$, because the two differ by a factor of two and the distinction decides both bounds below.

\emph{The declaration.} Suppose we declare $i \succ j$, so $\min_z \mathrm{lo}(z) > 0$, and suppose
the declaration is wrong, so $\Delta(z_0) \le 0$ for some $z_0 \in S$. Then
$\mathrm{lo}(z_0) > 0 \ge \Delta(z_0)$: that one setting's lower bound exceeded its own truth, which
happens with probability at most $\alpha_1$. Since $z_0$ is fixed by the response surface rather
than chosen from the data, no union over $|S|$ is needed. The false-declaration probability for either
pre-specified direction is at most $\alpha_1 = 0.025$. Because the rule may declare either ordering,
the probability of any false declaration is at most $2\alpha_1 = 0.05$ for a fixed finite $S$.

\emph{Coverage of the range} $[\min_z \Delta(z), \max_z \Delta(z)]$ is a different statement and
costs a factor of two on the tails rather than a factor of $|S|$:
$\min_z \mathrm{lo}(z) \le \mathrm{lo}(z_{\min}) \le \Delta(z_{\min})$ fails only if the interval at
the \emph{argmin} misses low, and symmetrically at the argmax, so the envelope covers the range with
probability at least $1 - 2\alpha_1 = 0.95$ however many settings $S$ holds. Both bounds are
conditional on the per-setting intervals being valid, which Sect.~\ref{subsec5_4} qualifies where a
variance term is imputed, and neither says anything about parameter values we did not sample --- the
limitation Sect.~\ref{subsec8_2} measures. Both
statements are conditional on the intervals at each $z$ being valid, which Sect.~\ref{subsec5_4}
qualifies where a variance term is imputed.
Throughout the empirical sections $S = F_{\mathrm{nom}}$ and the single setting is nominal, which is
the pairing the benchmark's own reporting implies.
A third rung, a Gaussian-process simultaneous band over the parameter axes, is introduced in
Sect.~\ref{sec8} for the case where the decisive parameter value may lie \emph{between} the settings
we sampled; the union bound's guarantee is of sampling-extremum type and does not cover that case.

\subsection{The error ledger}\label{subsec3_4}

Writing $[L_{ij}, U_{ij}]$ for the union-bound interval, the claim that $i$ is genuinely better in
practice requires three events to hold simultaneously: that the true parameter lies in the set the
bound ranges over, $E_\theta = \{z^* \in S\}$; that the evaluation intervals cover the simulated
differences uniformly over the set, $E_{MC}$; and that the sim-to-real bias is bounded,
$E_b = \{|\Delta^{\mathrm{real}}_{ij} - \Delta^{\mathrm{sim}}_{ij}(z^*)| \le b_{ij}\}$. On the
intersection every real difference lies in $[L_{ij} - b_{ij}, U_{ij} + b_{ij}]$, and \emph{if} the
three carry levels $1-\alpha$, $1-\gamma$ and $1-\beta$, a union bound gives
$1 - \alpha - \beta - \gamma$ without requiring independence.

The conditional is doing real work and we leave it explicit. This paper establishes none of the
three levels. $E_\theta$ would need a set with a coverage guarantee, and Sect.~\ref{subsec3_2}
explains why ours has none as computed; $E_{MC}$ is a simultaneous statement and we report per-pair
intervals; $E_b$ needs a sim-to-real bias bound we do not have. The ledger is therefore an accounting
of what a defensible claim would require, and a record of which parts we supply --- not a derivation
of a number attached to our results.

We are deliberate about which of these this paper establishes, and the accounting is less favourable
than the arithmetic above suggests.

$E_{MC}$ is what Sects.~\ref{sec5}--\ref{sec7} measure, and enumerating $\mathcal{C}$ removes one of
its two components exactly. Two caveats attach to it: our intervals are computed per policy pair, so
what we report is a collection of per-pair statements rather than the simultaneous event
$E_{MC}$ written above, and we label them as such throughout rather than applying a multiplicity
correction we have not performed.

$E_\theta$ is where the ledger is weakest, and there are two separate reasons, neither of which we
repair. The first is Sect.~\ref{subsec3_2}: $C_\alpha$ as computed selects its reference on the same
demonstrations it tests on, so it carries no $1-\alpha$ guarantee either, and the split-sample
variant retains a true minimiser $92.4\%$ of the time against a nominal $95\%$. The second is
Sect.~\ref{subsec4_4}: even a valid $C_\alpha$ would not transfer to $F_{\mathrm{nom}}$, the fibre
our verdicts actually range over, because $F_{\mathrm{nom}}$ is generated by a
point the calibration data reject --- nominal's ratio and delay are identified and lose to the grid
minimiser on both stacks. So for the sets we actually use we claim no coverage
probability for $E_\theta$ at all. What we can say is narrower and still useful: $F_{\mathrm{nom}}$
contains every setting the calibration evidence cannot distinguish from the benchmark's operating
point, so a verdict that survives the union over it is one the calibration evidence cannot overturn
\emph{without also moving off the operating point}. That is a conditional statement, and we keep it
conditional.

$E_b$ we do \textbf{not} establish: we have no per-pair sim-to-real bias bound, and
Sect.~\ref{subsec8_4}'s delimiting experiment is best understood as evidence that on at least one
pair $b_{ij}$ is large. A criterion built on a conditional $E_\theta$ and a per-pair $E_{MC}$ bounds
the ambiguity it can represent and no more, and that is the honest description of what the ladder in
Sect.~\ref{sec8} delivers.

%%=============================================================================%%
\section{The calibration protocol is structurally blind}\label{sec4}

\subsection{Setup}\label{subsec4_1}

Calibration of the kind used by these suites fits simulator parameters by replaying recorded real
demonstrations: the recorded end-effector targets are fed to the simulated controller, and the
parameters are chosen to minimise the deviation between the simulated and recorded trajectories. Let
$z$ collect the joint PD stiffness $k$, the damping $d$, the torque limit $\tau_{\max}$, the
execution delay, and the contact parameters (object friction and density). The demonstrations are
teleoperated, free-space and slow, which is the property that drives everything in this section: the
protocol does not choose its own excitation.

\paragraph{The objective, written out, because its units matter later.}
We inherit the reference implementation's system-identification objective. For demonstration $m$ of
$M$, with $T_m$ recorded steps, simulated end-effector position $p_t$ against recorded $\hat p_t$
and orientation $R_t$ against $\hat R_t$,
\begin{equation}\label{eq:loss}
L_m(z) = \frac{1}{T_m} \sum_{t=1}^{T_m} \big( e^{\mathrm{p}}_t + e^{\mathrm{r}}_t \big),
\end{equation}
where $e^{\mathrm{p}}_t = \lVert p_t - \hat p_t \rVert_2$ is a position error in \emph{metres},
$e^{\mathrm{r}}_t = \arcsin(\lVert R_t - \hat R_t\rVert_F/(2\sqrt2))$ is an orientation error in
\emph{radians}, and $L(z) = M^{-1}\sum_m L_m(z)$. Adding the two makes this a \textbf{composite}
objective: it adds a length to an
angle, weighting one radian as one metre. That weight is inherited, not derived, and the sum is
therefore not a distance, and in particular it must not be scaled by $1000$ and reported in
millimetres. Values of
$L$ and differences in $L$ are quoted below in \emph{composite units}, and wherever the physical size
of an effect matters we give the translation and rotation components separately. For orientation:
at the fitted setting $L = 0.0456$, of which the translation part is $19.1$~mm and the remainder is
the rotation term.

\subsection{A structural insensitivity, not a flat region of a noisy objective}\label{subsec4_2}

Two directions in $z$ leave $L$ unchanged to within magnitudes we quantify in
Sect.~\ref{subsec4_3}. We are careful about the modality: each rests on a mechanism with a stated
domain of validity, not on an exact identity of the full dynamics, and each is a property of the
\emph{protocol} rather than of the numerical optimiser --- which is the distinction that matters,
because no improvement in fitting removes it.

\paragraph{(i) Common scaling of the PD gains.}
Under free-space replay at demonstrated velocities, the joint response to a target is dominated by a
first-order lag with time constant $d/k$: the inertial, gravitational and Coriolis terms are small
relative to the PD action at these speeds, and to that order scaling $(k, d) \mapsto (ck, cd)$ leaves
the time constant, hence the trajectory, unchanged. \emph{This is an approximation, and we state its
limits.} Under the full rigid-body dynamics a common gain rescaling does not preserve the trajectory
exactly --- the inertial term does not scale with the gains --- so the invariance is one of
\emph{degree}, holding while the neglected terms stay small compared with the calibration protocol's
own resolution. What makes it a structural rather than an incidental property is that the smallness
follows from how the demonstrations were collected --- teleoperated, slow, free-space --- and not
from a fortunate parameter value. Sect.~\ref{subsec4_3} measures the residual: over
$\times \tfrac14$ to $\times 4$ it is at most $3.1\times10^{-6}$ in the loss, against $1.36\times10^{-3}$ between
two implementations of the identical nominal dynamics. Outside that range, and at speeds that excite
the inertial terms, we make no claim; a faster demonstration set would break it, which the remark
below treats as the point rather than as a caveat.

\paragraph{(ii) The torque limit.}
Wherever the commanded torque stays strictly below $\tau_{\max}$ the saturation is inactive and the
trajectory does not depend on its value at all --- not to a tolerance, but because the limit never
enters the computation. This one is exact where it applies, and what bounds it is the domain rather
than the mechanism: the argument holds for every demonstration that does not approach the limit, and
says nothing about any that does. It therefore supports insensitivity over an \emph{interval} of
$\tau_{\max}$ --- down to the largest commanded torque the demonstrations actually produce --- and
not down to arbitrarily small limits, which would eventually clip the demonstrated motion. We
measured it at $\tau_{\max}/2$, where 97 of our 98 demonstrations fall in the inactive case
(Sect.~\ref{subsec4_3}); the 98th is evidence that the interval's edge is not far below that, and we
claim nothing past the value we tested.

Both invariances are properties of the protocol's \emph{excitation range}, not of the simulator. A
faster or more contact-rich demonstration set would excite gravitational and Coriolis terms that
break the first, and would drive the commanded torque into saturation, breaking the second. That is
the point rather than a caveat: the blindness is inherited from demonstrations that were collected
for teleoperation rather than for identification, and it would shrink under an excitation designed
to identify. What cannot shrink it is more of the same data, or a better fit to it --- which is the
distinction between coverage and precision that Appendix~\ref{secA} makes formal.

The ratio direction $d/k$ and the execution delay \emph{are} identified, and both produce the
expected bowl-shaped loss curves; the claim is specific to the two directions above.

\subsection{Empirical confirmation on two independent simulator stacks}\label{subsec4_3}

We replayed 98 BridgeData V2 demonstrations \citep{walke2023bridgedata} on (a) the ManiSkill3 port
\citep{tao2024maniskill3} and (b) the original reference stack (ManiSkill2
\citep{gu2023maniskill2} with SAPIEN \citep{xiang2020sapien}), the latter running headless through a
purpose-built Vulkan layer (Appendix~\ref{secB}).

\paragraph{What these 98 demonstrations are, and what they are not.}
This bears on how strongly the rest of this section can be read, so we state it before the results
rather than in a limitation. The demonstrations are read from four fixed shards of the public
BridgeData V2 RLDS release by bounded HTTP range request, with CRC32C and SHA-256 verified per
record, and the manifest records the byte offset of every one. They are therefore \emph{the same
public dataset} the reference calibration draws on, and the selection is reproducible to the byte.
They are \textbf{not} verified to be the same trajectories the reference parameters were fitted on.
The reference preparation script selects its demonstrations by TFDS iteration index; we read records
positionally from shards. Record $k$ of shard 0 very probably is TFDS iteration $k$ under an
unshuffled read, but the format does not guarantee it, our loader records the equality as an
assumption rather than a check, and we could not locate a fitting log or trajectory-identity map to
settle it. Nor is our search domain the reference's: we enumerate a 50-point grid that may be wider
than the one the original fit searched.

Two consequences, and we hold to them throughout. What this section establishes is that
\emph{applying the reference replay-calibration protocol to 98 demonstrations from the same public
dataset, over our stated grid, the shipped controller setting is disfavoured on both simulator
stacks.} What it does not establish is the stronger historical claim that the original fitting
evidence, on the original sample and under the original search domain, would have rejected the
setting that fit produced. Replicating on two stacks and splitting the sample address precision and
selection; neither addresses sample identity, and a conclusion drawn on 98 demonstrations would not
transfer automatically to a 12-demonstration subset even if that subset were shown to be contained
in ours. The consequence this paper actually needs is the weaker one: the benchmark's operating point
is not the one \emph{a} faithful application of its own stated protocol prefers, which is enough to
make the question of Sect.~\ref{subsec7_7} --- does the ranking move with the operating point? ---
worth measuring.

\begin{table*}[t]
\caption{Replay identifiability on two independent simulator stacks. 98 BridgeData V2
demonstrations. The two insensitive directions are flat to sub-millimetre and to bitwise on 97 of 98
demonstrations respectively; the identified directions produce the expected loss increase. The scale
to judge these against is the $1.36\times10^{-3}$ composite disagreement between these two stacks at
\emph{identical} nominal parameters (Sect.~\ref{subsec4_4}), whose translation part is $0.47$~mm.
Sect.~\ref{subsec4_4} states what that magnitude is and is not: an engineering reference with an
external basis, not a resolution limit below which an effect becomes unattributable.}\label{tab1}
\begin{tabular}{@{}llll@{}}
\toprule
Direction & ManiSkill3 & Original stack & Verdict \\
\midrule
Common scale $c \in \{0.25, 0.5, 2, 4\}$ & max traj.\ diff.\ 0.141 mm & 0.207 mm & invariant \\
Torque limit $\times 0.5$ & bitwise 97/98, worst $0.149\,\mu$m & 97/98, worst $0.050$ mm & invisible \\
Ratio $d/k \times \{0.5, 2\}$ & $+0.0089$ / $+0.0064$ & same sign and order & identified \\
Execution delay 1 step & clearly identified & clearly identified & identified \\
\bottomrule
\end{tabular}
\end{table*}

All 14 iso-scale groups have maximum trajectory differences below $1$~mm; we quote that threshold
as a round descriptive figure, not as a tolerance anyone has defended. The paired \emph{loss} increases for
the ratio direction are $+0.0089\,[+0.0073, +0.0105]$ and $+0.0064\,[+0.0055, +0.0072]$ in composite
units --- three orders of magnitude above the $3.1\times10^{-6}$ the common scale produces over a
larger range, which is the contrast this section exists to establish. Figure~\ref{fig1} shows the
bowl-shaped loss along the identified directions against the flat loss along the insensitive ones.

Three quantities appear in this section and they are not interchangeable, so we name them once. The
$0.141$ and $0.207$~mm in Table~\ref{tab1} are maximum \emph{end-effector position} differences
between parameter settings, in metres. The loss increases just quoted are differences in the
\emph{composite} objective of Eq.~\eqref{eq:loss}, which is not a length. And neither is a residual
of the fit: the fit's own residual at the fitted setting is $L = 0.0456$ composite, of which
$19.1$~mm is translation. Sect.~\ref{subsec4_4} judges parameter effects against a between-stack
reference magnitude in the same composite units. Each is quoted in its own terms throughout and none
is a substitute for another.

Halving the torque limit leaves the replayed trajectory \emph{bitwise} identical on 97 of the 98
demonstrations, on both stacks independently, consistent with what Sect.~\ref{subsec4_2}(ii) predicts
for a saturation that never activates. The exception is instructive and we report it rather than
round it away: it is the same demonstration on both stacks, and there the trajectories differ by
$0.149\,\mu$m on ManiSkill3 and $0.050$~mm on the original stack.

We are deliberately cautious about why. That the same demonstration is the exception on both stacks
is suggestive of a property of that demonstration rather than of per-stack numerical noise, but it is
not decisive: two implementations of the same equations replaying the same trajectory can share a
sensitivity, and the common cause could be the trajectory rather than the torque limit. The
commanded torques were not recorded in the replay archive, so \textbf{we cannot confirm the
saturation mechanism from the stored data and do not claim to}; it remains the natural explanation
and an untested one, and re-running the replay with torques logged would settle it cheaply. What the
data do support without an assumption is the magnitude: the trajectory is unchanged on 97 of 98
demonstrations and changes by $0.050$~mm on the remaining one, which is a ninth of the $0.47$~mm of
translation inside the composite disagreement between the two stacks at identical parameters, and
well inside any tolerance the fit could resolve
(\texttt{scripts/audit\_exact\_invariance.py}).

\begin{figure*}[t]
\centering
\includegraphics[width=\textwidth]{Fig1}
\caption{Replay identifiability. Loss along the identified directions (the gain ratio $d/k$ and the
execution delay) is bowl-shaped with a well-defined minimum; loss along the invariant directions
(common gain scaling, torque limit) is flat --- for the torque limit, bitwise flat on 97 of the 98
demonstrations and within $0.050$~mm on the remaining one. Panels share the vertical
axis.}\label{fig1}
\end{figure*}

\subsection{The compatible set, and the fact that it excludes the benchmark's own operating
point}\label{subsec4_4}

Inverting the calibration test requires a search domain, a loss, a reference and a threshold, and the
answer depends on all four, so we state them. The domain is the 50-point grid of
$(\text{stiffness scale}, \text{damping scale}, \text{execution delay})$ that the replay sweep
enumerates; the loss is the per-demonstration composite objective of Eq.~\eqref{eq:loss}, which sums
a position error in metres and an orientation error in radians; the reference is the
grid's loss minimiser; and a candidate is retained when the one-sided $95\%$ lower bound of its
paired mean loss increase over that minimiser does not exceed zero, which is the rule
Appendix~\ref{secA} states. Computed that way on the 98 common demonstrations
(\texttt{scripts/rebuild\_compatible\_set.py}), the loss has a sharply two-part structure.

\paragraph{The loss depends on $(d/k, \text{delay})$ and on nothing else.}
Grouping the 50 grid points by ratio and delay gives 26 groups, of which 14 contain two or more
points that differ only in the common scale. Across those 14 --- spanning seven distinct ratios from
$0.5$ to $2.0$ --- the mean loss moves by at most $7.2\times10^{-6}$ within a group, identically on
both stacks. \emph{Between} groups it ranges over $6.0\times10^{-2}$, which is $8{,}354$ times the
largest within-group spread on ManiSkill3 and $8{,}767$ times on the original stack. Extending to
the sixteenfold range the policy sweeps actually use ($\times\tfrac14$ to $\times 4$, measured on
the original stack), the largest mean paired loss difference from nominal is $3.1\times10^{-6}$.
All loss figures in this section are in the composite units of Eq.~\eqref{eq:loss}, not in metres.

This is the insensitivity of Sect.~\ref{subsec4_2}(i) measured on the calibration objective itself
rather than on a trajectory, and it is the sharpest form in which we can state it. Note what it is
and is not: a property verified at seven ratios, not proved at all of them. Ratio $0.25$ has exactly
one grid point per delay value, so the common scale is not varied there at all --- which matters in
Sect.~\ref{subsec7_7}, because $0.25$ is the ratio the calibration data prefer.

\paragraph{The threshold in Eq.~\eqref{eq:cset} degenerates on a deterministic simulator.}
Before reading the set off, the rule has to be examined, because at threshold zero it does something
indefensible. Replaying demonstration $m$ at parameter $z$ is deterministic --- the same loss every
time --- so the only randomness is which demonstrations were drawn, and with 98 paired
demonstrations the bootstrap resolves mean differences of order $10^{-7}$ in composite loss. At
threshold zero the rule therefore \emph{rejects} $\times 2$ and $\times 4$ common rescalings whose
mean loss differs from nominal by $4.9\times10^{-7}$ and $7.4\times10^{-7}$ --- some two to three
thousand times below $\tau$ --- and as the demonstration count grows the retained
set shrinks to the single loss minimiser whatever the physics. Statistical significance is the wrong
question to ask of a difference that carries no noise; magnitude is the right one.

\paragraph{A reference magnitude whose basis is external to the conclusion.}
Since the threshold cannot be zero, something has to set the scale, and we would rather it were not
us. The two stacks implement the \emph{same} nominal dynamics, and on the same 98 demonstrations
their losses differ by $1.36\times10^{-3}$ in composite units, $95\%$ CI
$[0.61, 2.10]\times10^{-3}$, decomposing into $0.47$~mm of translation and $0.88$~mrad of rotation;
the mean absolute per-demonstration disagreement is $2.76\times10^{-3}$. Two independent
implementations of one set of equations differ by that much, so we use it as a reference magnitude
$\tau = 1.36\times10^{-3}$, reading the set at $\tau$ as well as at zero and reporting
$2.10\times10^{-3}$ as a sensitivity.

We state its status carefully, because it is easy to overstate. This is an
\emph{engineering reference} with an external basis, not a resolution limit: a same-stack paired
parameter effect can be both smaller than a fixed cross-stack offset and perfectly well resolved,
since the two quantities have different sources of variation. What $\tau$ buys is a scale that was
not chosen to produce an answer. It does not license the claim that anything below it is
unattributable. Nor is the choice among the signed mean, the mean absolute difference and a quantile
forced --- they differ here by a factor of two --- and we report which we used.

\paragraph{The insensitive directions survive every rule; the nominal setting does not clearly
survive any.}
Against $\tau$, the iso-scale family ($\le 3.1\times10^{-6}$) is retained by a factor of $440$ and
the halved torque limit ($\le 1.8\times10^{-8}$) by a factor of $75{,}000$. No conclusion in
Sects.~\ref{subsec4_2} and~\ref{subsec4_3} depends on the threshold.

Nominal is the opposite case. The grid's loss minimiser is not the simulator's shipped controller:
on \emph{both} stacks independently it lies at one step of execution delay and a gain ratio of
$0.25$ to $0.36$, while nominal sits at ratio $1$ with no delay and ranks fifth of 26 groups.
Nominal's paired mean loss increase over the minimiser is $+2.90\times10^{-3}$ with a $95\%$ lower
bound of $+2.25\times10^{-3}$ on ManiSkill3, and $+2.26\times10^{-3}$ with $+1.49\times10^{-3}$ on
the original stack; the ManiSkill3 figure decomposes into $1.06$~mm of translation and $1.84$~mrad
of rotation. Both bounds exceed zero and both exceed $\tau$; at the $2.10\times10^{-3}$ sensitivity
nominal is still out on ManiSkill3 and comes in on the original stack. So its membership is
\textbf{borderline}, and the exact grid point depends on the objective's rotation weight: under
weights from $0$ to $5$ radians per metre the ManiSkill3 minimiser is unchanged while the
original-stack minimiser moves among three adjacent points. \emph{What survives every weight is the
qualitative statement} --- both stacks prefer one step of delay and a ratio well below $1$, and
neither prefers nominal. We report that rather than a single privileged setting, because ``the point
the data identify'' is only defined once the objective is.

\paragraph{The exclusion survives a procedure whose reference is selected independently.}
The rule above selects its reference from the same demonstrations it tests on, which
Sect.~\ref{subsec3_2} shows invalidates the coverage reading, so the exclusion has to be checked
against something better before any weight is put on it. Selecting the reference on one half of the
98 demonstrations and testing on the other gives independent selection and inference. Nominal is
still excluded: on ManiSkill3 its paired mean loss increase over the half-sample reference is
$+2.80\times10^{-3}$ composite with a one-sided lower bound of $+1.90\times10^{-3}$, and on the
original stack $+2.49$ and $+1.43\times10^{-3}$. Across five random splits the lower bounds run from
$+1.02$ to $+3.03\times10^{-3}$ and nominal is
excluded in all ten stack--split combinations, with the minimiser at delay $1$ and ratio $0.25$ or
$0.35$ every time. The conclusion of this subsection does not depend on the invalid step, which is
why we keep it.

Two consequences follow, and we separate them carefully. Our earlier analysis fixed nominal \emph{as}
the test reference, which makes its own loss increase zero by construction and so could never reject
it; the issue was invisible for that reason, not because the data spoke in its favour.

First, the invariance results of Sects.~\ref{subsec4_2} and~\ref{subsec4_3} are untouched and in fact
strengthened, for the reason just given: they are statements about directions along which the loss
does not move, they hold under every threshold we examined, and nothing about them depends on where
the minimum sits.

Second, the object our verdicts range over is \emph{not} $C_\alpha$, and we stop calling it that. All
of our policy sweeps are anchored at nominal: the conditions we evaluate are nominal together with
the settings that leave the replay loss unchanged relative to it --- common rescalings
$\times 0.25$ and $\times 4$, the halved torque limit, and contact parameters the free-space replay
never sees at all. That collection is the \emph{calibration-invisible fibre through the benchmark's
operating point},
\begin{equation}\label{eq:cset2}
\begin{split}
F_{\mathrm{nom}} = {}
  &\{d/k = 1,\ \text{delay}\ 0\}
   \times \{\text{scale} \in [\tfrac14, 4]\} \\[-2pt]
  &\times \{\tau_{\max}/\tau_{\max}^{\mathrm{nom}} \in [\tfrac12, 1]\} \\[-2pt]
  &\times (\text{contact axes, unobserved}),
\end{split}
\end{equation}
where the scale and torque intervals are the ranges over which we \emph{measured} the loss to be
flat, not the ranges over which we assume it. This matters: the invariance argument of
Sect.~\ref{subsec4_2} is a statement about a mechanism, and a mechanism has a domain. A common gain
rescaling leaves the free-space response unchanged to the order the demonstrations excite, which our
measurements confirm over $\times \tfrac14$ to $\times 4$ and which we do not claim beyond it; an
inactive saturation limit is invisible only while it stays inactive, which bounds the torque
statement below rather than at arbitrarily small $\tau_{\max}$. The contact axes are different in
kind --- unobserved rather than measured flat --- and remain unbounded, which is a statement about
the protocol's coverage.

$F_{\mathrm{nom}}$ is the right object for the question we actually ask --- \emph{given that the
benchmark runs at nominal, what does the calibration evidence determine about the numbers it
publishes?} --- because every published rate is computed at nominal. It is not the right object for
the different question of what the calibration evidence determines about the simulator, and we do not
use it for that. $F_{\mathrm{nom}}$ and $C_\alpha$ overlap only in the axes along which the loss is
flat; they disagree on the ratio and the delay, which the data identify and on which nominal loses.

This raises a question we could evade but should not: if the benchmark's operating point is one the
calibration data reject, does the ranking computed there survive a move to the point they prefer? We
answer it by measurement rather than argument. Sect.~\ref{subsec7_7} re-runs the complete
64-configuration eggplant census at $(k{\times}2,\ d{\times}0.5,\ \text{delay}\ 1)$, on the same
configurations and the same policy seeds, so that the two operating points differ in nothing else.
The answer splits by criterion. The point-calibration verdict does not survive: it moves from
abstention to a declaration. The set-valued verdict does, abstaining at both points --- which is
what a criterion that reports the range of answers the evidence permits is for.

\textbf{The question of this paper is therefore what a policy ranking computed over
$F_{\mathrm{nom}}$ is worth}, with the separate finding that the operating point generating that
fibre is one our re-application of the calibration protocol disfavours --- in the bounded sense
Sect.~\ref{subsec4_3} states --- and with Sect.~\ref{subsec7_7} measuring what that costs.

%%=============================================================================%%
\section{What a finite benchmark measures}\label{sec5}

\subsection{The benchmark is a finite population of configurations}\label{subsec5_1}

Both widely used ports map the episode index onto a fixed grid of initial states,
\begin{align}
\mathrm{pos}  &= \big\lfloor (\mathrm{episode\_id} \bmod (n_{xy} n_{quat})) / n_{quat} \big\rfloor, \\
\mathrm{quat} &= \mathrm{episode\_id} \bmod n_{quat},
\end{align}
so the number of distinct initial configurations is $n_{xy} \cdot n_{quat}$, and any episode index
beyond that repeats an earlier scene exactly. Object scale is drawn from the asset table, but every
object has a single admissible scale, so it adds no variation.

\begin{table*}[t]
\caption{Distinct initial configurations per task in the two ports. The eggplant task differs in both
the number of orientations and the quaternion values, so an episode index does not denote the same
scene across ports.}\label{tab2}
\begin{tabular}{@{}lrr@{}}
\toprule
Task & Original stack & ManiSkill3 port \\
\midrule
Put eggplant in basket & $8 \times 3 = \mathbf{24}$ & $8 \times 8 = \mathbf{64}$ \\
Put spoon on towel & $12 \times 2 = 24$ & $12 \times 2 = 24$ \\
Put carrot on plate & $12 \times 2 = 24$ & $12 \times 2 = 24$ \\
Stack green on yellow cube & $24 \times 1 = 24$ & $24 \times 1 = 24$ \\
\bottomrule
\end{tabular}
\end{table*}

Three consequences follow, and all three are visible in published practice.

\begin{enumerate}
\item \textbf{Beyond the grid, episodes are exact repeats.} A ``96-episode'' evaluation of the
eggplant task on the ManiSkill3 port is 64 distinct scenes, 32 of which are visited twice.
\item \textbf{The two ports do not evaluate the same benchmark.} For the eggplant task they differ
in the number of orientations (3 versus 8) and in the quaternion values themselves: the best matching
pair has $|\langle q_1, q_2 \rangle| = 0.54$. Per-episode agreement between the ports is therefore
meaningless for that task. Spoon and carrot share their grids exactly --- and their cross-port
per-episode agreement is 0.71--0.96, against chance level for eggplant. That contrast is the control
for this diagnosis.
\item \textbf{One scene is worth $1/24$ or $1/64$.} On a 24-configuration task, any reported
difference below 0.042 is less than one scene.
\end{enumerate}

\subsection{The official protocol is already a census}\label{subsec5_2}

The reference evaluation script runs each bridge task over episode range $0$ to $24$ under three
fixed policy seeds: the \textbf{complete 24-configuration census repeated under three fixed policy
seeds}, 72 episodes, which is exactly the denominator of the published table (for instance $31/72 =
0.431$ and $41/72 = 0.569$ for the two Octo variants on the eggplant task). The ManiSkill3 port's
evaluation script instead defaults to 100 episodes indexed as $\mathrm{seed} + i$, aligning with
neither the grid nor the seed repetition. The community's own reference protocol already does the
right thing; the drift is in what downstream work does with it.

\subsection{What counts as a replicate depends on the policy}\label{subsec5_3}

A policy seed is a replication axis only if the policy consumes it. Octo's diffusion action head
draws a pseudorandom key that enters sampling, so changing the seed produces a new sample and
configuration $\times$ seed is the legitimate replication axis. OpenVLA decodes greedily with
sampling disabled, so the seed is unused and changing it produces nothing new; configuration is the
only axis. Empirically, two ``seed sets'' of the deterministic policy are record-identical in 48 of
48 episodes. For such a policy, ``replicated across seed sets'' is vacuous, and the configuration
grid --- which can be enumerated --- is the only sampling axis. This is the error behind the second
retraction in Sect.~\ref{subsec1_2}.

\subsection{Two estimands, two intervals}\label{subsec5_4}

Two different quantities are both legitimate, and they carry different uncertainty.

\emph{(i) The benchmark value} --- the number this benchmark reports. After a census, configuration
sampling error is zero and only policy noise and run-to-run numerical noise remain, so
\begin{equation}\label{eq:var}
\mathrm{Var}(\hat\Delta) = N^{-2} \sum_{c} \big[ s^2_A(c)/S_c + s^2_B(c)/R_c \big],
\end{equation}
with $S_c$ and $R_c$ the number of runs of each policy at configuration $c$. ``Run'' means one
episode, and the count is pooled across seed-set directories: a configuration that received six
episodes contributes $S_c = 6$, even where those episodes came from fewer directories. That choice
matters --- an implementation that instead averaged within a directory and took each directory mean
as the observation moves two verdicts --- and Sect.~\ref{subsec7_2} reports both alongside the
reason for preferring this one.

\emph{(ii) The generalisation value} --- performance on comparable scenes. Here configuration
sampling error is real and the appropriate estimator is a paired bootstrap \emph{clustered by
configuration}.

One term in Eq.~\eqref{eq:var} has to be imputed and we state the rule rather than leave it in the
code. A configuration with a single run supplies no within-configuration variance, so $s^2(c)$ is
unavailable there. Where the cell has other configurations with two or more runs we substitute the
mean of the measured $s^2$; where \emph{no} configuration has a replicate we fall back to the
binomial $\bar p(1 - \bar p)$; and for a policy that decodes deterministically, whose re-runs are
record-identical, we set the term to zero and label the cell as carrying no measurement of run-to-run
drift. Of the 18 rows in the release's core table, 15 are fully empirical on both sides, 2 use the
deterministic-zero term for the OpenVLA side, and 2 use the binomial fallback. Every row prints its
model, so a reader can see which interval rests on an imputation; the imputation is a pooled
variance, which is an assumption of homogeneity across configurations that we do not test and that
would be optimistic if run-to-run noise were concentrated in a few scenes.

Three further implementation points matter and are easy to get wrong. With a single run per configuration,
estimator (i) degenerates to the ordinary binomial interval: \textbf{a census does not narrow an
interval; replicates do.} Enumerate the grid first, which removes double counting, then add runs,
which removes policy noise. Record-identical re-runs of a deterministic policy must be merged into
one observation, or the variance is halved for free. And estimator (ii) must cluster by
configuration; in our data this widens intervals by a median factor of 1.04 on the 64-configuration
task and up to 1.33 on a 24-configuration task run for 96 episodes, and changes no verdict --- but
the effective sample size must be reported honestly.

\subsection{The evaluation protocol and the implementation build are themselves invisible
parameters}\label{subsec5_5}

Holding the configuration grid fixed, we separated three further sources.

\begin{remark}[Scope of the quantities in this paper]\label{rem:scope}
Unless a number is explicitly labelled otherwise, every effect size and interval we report is
census-scope at five seed sets where five exist: the task's configuration grid is fully enumerated,
so configuration sampling error is exactly zero, and the policy-noise term is estimated across as
many independent seed sets as that task has. Five is the count for the eggplant census; the spoon
and carrot sweeps have two, and Appendix~\ref{secC} gives the per-cell figures. Where a quantity is also available at the three-seed budget the reference protocol uses, we
give both and say which is which, because the difference between them is itself one of the paper's
results (Sect.~\ref{subsec7_2}).
\end{remark}

\begin{table*}[t]
\caption{Sources of uncertainty that persist after the configuration grid is enumerated. The three
magnitudes are \emph{different kinds of quantity} --- a standard deviation across seed sets, a paired
mean shift between builds, and the spread of one policy's rate across four settings --- so the column
reports each in its own terms and they should not be ranked against one another as though they shared
a scale. Every entry is a quantity in units of $\Delta$, the policy difference, and not a change in a
single policy's success rate --- the two are easy to interchange and are not the same magnitude.
The scopes differ per row. The seed row
is census-scope at $S = 5$ on eggplant, median over policy pairs, computed on the episode ids common
to all five seed sets; \texttt{analyze\_seed\_noise.py} reports $0.068$ for the same statistic on its
own per-pair id intersection, and the third-decimal difference is that restriction rather than a
rounding of one number. The build row is a single run of
each of two server builds on the eggplant census, so it carries no seed replication and no $S$; it
reproduces at the episode level on 0.73--0.88 of episodes. The texture row is one run per
configuration over the 300-configuration fractal census, also without seed replication; the four
variants give $\Delta = +0.213, +0.120, +0.200, +0.053$ there, while one policy's \emph{rate} over
the same four spans $0.400$ (0.000, 0.280, 0.013, 0.400) --- a larger number that belongs to a
different quantity and is quoted in the text, not here. Seed and build rows are paired on the
episode ids the two directories share
(\texttt{analyze\_platform\_drift\_paired.py}).}\label{tab3}
\begin{tabular}{@{}llll@{}}
\toprule
Source & Isolation & Quantity & Magnitude \\
\midrule
Policy seed & different seeds, same build & sd of $\Delta$ over seed sets & $0.067$ \\
Implementation build & same seeds, newer build & shift in $\Delta$ & $-0.109$ \\
Robot texture variant & \texttt{urdf\_version} only & span of $\Delta$ over 4 variants & $0.160$ \\
\bottomrule
\end{tabular}
\end{table*}

The third row needs unpacking. \texttt{urdf\_version} selects among four recolourings of the robot
model. It changes no physical property --- not a mass, not a friction coefficient, not a joint limit
--- and the published figure for this task is the mean over all four. On the fractal coke-can census
the four variants give one policy 0.000, 0.280, 0.013 and 0.400, so evaluating on a single variant
displaces the reported number by $-0.173$ to $+0.227$ relative to the four-variant mean. For
comparison, the can's physical orientation, which does change the dynamics, spans only 0.140.

This axis is invisible to replay calibration for the same reason the torque limit is
(Sect.~\ref{subsec4_2}): it never enters the demonstration trajectory. Unlike the torque limit it is
not even a physical parameter, which makes it the paper's sharpest illustration that
``calibration-invisible'' is a property of the protocol's coverage rather than of physics. That a
visual variant should affect a vision-conditioned policy is not itself surprising --- it is the
premise of domain randomisation \citep{tobin2017domain, peng2018sim}.

Nor is the averaging itself undocumented: the suite describes its visual-matching design and reports
the variant aggregation as part of the protocol \citep{li2024simpler}. Our contribution here is not
the discovery of a hidden step but a statement about the \emph{estimand} that step defines, and we
want to state it precisely, because the loose version is wrong. The average over four variants is
still a success rate --- a perfectly well-defined one. What it is not is the rate \emph{conditional
on} a variant. The published number is the rate \emph{marginal over} a uniform mixture of the four,
and on this task the four conditionals span 0.400.

Two consequences follow, and neither is a variance claim. First, a reader comparing a published rate
against a number obtained on one variant --- the natural thing to do when reproducing, since a
single run fixes a single \texttt{urdf\_version} --- is comparing a marginal against a conditional.
The discrepancy that results is $-0.173$ to $+0.227$ here with nothing having gone wrong in either
computation. Second, and this is where it is easy to slip: under a complete enumeration the marginal
carries \emph{no} sampling error from the variant axis at all, because every variant is run. The
0.400 spread is therefore not uncertainty in the published number; it is the distance between that
number and the four quantities a single-variant evaluation would produce. Treating the four as
replicates of one quantity --- and reporting a standard error over them --- would misdescribe both.
What we add is the magnitude and this bookkeeping, alongside the suite's own observation that visual
matching is what its simulated rates depend on.

Whether the axis can flip a \emph{ranking} rather than widen one policy's interval is a separate
question, and Sect.~\ref{subsec8_4} answers it on a second policy over the same grid. It can move a
verdict to abstention, but it does not reverse an ordering: all twelve texture $\times$ condition
cells there keep the same sign.

Both the seed and build rows are computed on matched episode scopes. That is not a formality: the
directories differ in size --- the pre-fix build and one seed set hold 96 episodes against 64
elsewhere --- and because episode ids beyond the grid wrap onto it, averaging each directory over its
own contents compares different effective scopes, which changes the answer by roughly a factor of
two. We report the paired figures throughout and the unpaired ones nowhere, and
\texttt{analyze\_platform\_drift\_paired.py} prints both so the choice is inspectable. We do not,
however, use the comparison to rank the build effect against policy-seed noise: a shift in $\Delta$
and a standard deviation of $\Delta$ are different quantities, and which is ``larger'' is not a
well-posed question.

Two further practical consequences follow. First, one seed set plus a per-episode bootstrap declares
rankings that do not hold: in our own data $\Delta(\text{small}-\text{base}) = +0.250\,[+0.08,
+0.32]$ declared one Octo variant better, and three seed sets from a single build give $+0.141 /
-0.031 / -0.047$. We retract that verdict. Second, replicates must never be pooled across builds;
doing so inflated the measured seed standard deviation from 0.055 to 0.093, so about half of the
apparent seed noise was build drift.

\begin{figure*}[t]
\centering
\includegraphics[width=\textwidth]{Fig2}
\caption{The uncertainty budget after the configuration grid is enumerated. The point of the figure
is which terms \emph{respond to evaluation budget}: the three evaluation-noise bars shrink as
$1/\sqrt{S}$, while the implementation build, the calibration-invisible torque limit and the robot
texture variant do not shrink at all. Every bar is a quantity in units of $\Delta$, the policy
difference, but they are three different \emph{kinds} of such quantity --- a 95\% half-width, a
shift, and a span across four settings, as each label states. They share an axis because they share
a unit, not because they are interchangeable, and their lengths should not be ranked against one
another. Scopes differ per bar and are given in Table~\ref{tab3}. Every value is computed from the
episode records.}\label{fig2}
\end{figure*}

\subsection{Reporting requirements}\label{subsec5_6}

Every number in this paper carries: configurations covered / total; independent runs per
configuration, and whether they are genuinely independent for that policy; the variance model
(empirical, binomial fallback, or unmeasured); which estimand it targets; and the port and build it
was collected on. We log the last of these per run in an append-only record, having discovered that
a per-directory metadata file is overwritten by the next run and makes the provenance of
already-collected episodes unrecoverable.

%%=============================================================================%%
\section{Reproducing the published protocol}\label{sec6}

\subsection{The protocol, followed exactly}\label{subsec6_1}

The reference implementation's bridge evaluation script evaluates each task over episode range $0$ to
$24$ under three values of \texttt{init\_rng}: a complete enumeration of the task's 24 initial
configurations, repeated three times, for the 72 episodes that form the denominator of the published
rates.

How the seed is used inside a sweep is easy to get wrong, and we got it wrong first. The reference
wrapper seeds \emph{once}, in its constructor, and then lets one pseudo-random stream advance through
every step of every episode of the run; its \texttt{reset()} restores the task, the image history,
the action ensembler and the sticky-gripper state, and never touches the key. So a run's 72 episodes
carry 72 distinct noise realisations, and \texttt{init\_rng} indexes a whole stream rather than a
per-episode draw. Our first reproduction instead sent the seed on every reset, and our inference
server re-seeded on receipt --- with the consequence that all 24 configurations within a run replayed
one identical noise realisation. We report both, because the difference is not cosmetic: it changes
the agreement with the published table, the across-run spread, and one of our own conclusions
(Sect.~\ref{subsec6_3}). All numbers below are from the reference lifecycle unless labelled
otherwise.

We reproduced this on the original reference stack, running headless through the Vulkan compatibility
layer of Appendix~\ref{secB}, with the policies served by our own inference server rather than
in-process. Published values are read from the pinned metrics table distributed with the benchmark.
\texttt{scripts/queue\_ms2\_official\_stream.py} runs the reference lifecycle and asserts, on every
episode, that the server applied the lifecycle it asked for; \texttt{analyze\_official\_protocol.py
--compare} prints the two side by side.

\subsection{The reproduction agrees, and the differences are one-signed}\label{subsec6_2}

\begin{table*}[t]
\caption{Our 72-episode reproduction against the published values, under the reference RNG
lifecycle. The rightmost column gives the same pooled rate from our earlier per-episode-re-seeding
run, for comparison. Under the reference lifecycle the mean observed difference from the published
table across these eight cells is $-0.038$ (mean absolute 0.038, largest 0.111), and six of the
eight are below the published value with two at zero; under the re-seeding lifecycle the mean is
$-0.005$ and the mean absolute 0.026. The eight cells are four tasks $\times$ two policies, sharing
tasks, policies and seeds, so they are not eight independent replications and we quote the spread
across them descriptively rather than as a confidence statement about a single platform
parameter. Note also that the eggplant octo-base cell agrees to $0.4306$ against a published
$0.431$, which is a rounding difference and not evidence of a downward shift.}\label{tab4}
\begin{tabular}{@{}llrrrrrrr@{}}
\toprule
 & & \multicolumn{4}{c}{reference lifecycle} & & & re-seeding \\
\cmidrule(lr){3-6}
Task & Policy & seed 0 & seed 2 & seed 4 & pooled & published & diff. & pooled \\
\midrule
eggplant & octo-small & 0.542 & 0.500 & 0.542 & 0.528 & 0.569 & $-0.041$ & 0.556 \\
eggplant & octo-base & 0.458 & 0.417 & 0.417 & 0.431 & 0.431 & $-0.000$ & 0.472 \\
spoon & octo-small & 0.375 & 0.417 & 0.333 & 0.375 & 0.472 & $-0.097$ & 0.431 \\
spoon & octo-base & 0.042 & 0.000 & 0.000 & 0.014 & 0.125 & $\mathbf{-0.111}$ & 0.083 \\
carrot & octo-small & 0.042 & 0.042 & 0.167 & 0.083 & 0.097 & $-0.014$ & 0.069 \\
carrot & octo-base & 0.042 & 0.083 & 0.042 & 0.056 & 0.083 & $-0.027$ & 0.125 \\
stack cube & octo-small & 0.000 & 0.042 & 0.042 & 0.028 & 0.042 & $-0.014$ & 0.042 \\
stack cube & octo-base & 0.000 & 0.000 & 0.000 & 0.000 & 0.000 & $+0.000$ & 0.000 \\
\bottomrule
\end{tabular}
\end{table*}

Two things in Table~\ref{tab4} matter, and they pull in opposite directions.

The magnitudes are small. The policy-seed noise we measure on the 64-configuration eggplant census is
$\pm 0.132$ at half-width for a single seed set and $\pm 0.076$ for three
(\texttt{analyze\_seed\_noise.py}), and every cell's discrepancy is inside that band. Those two
figures come from the spread of $\hat\Delta$ \emph{across whole seed sets} --- median
$\mathrm{sd} = 0.068$ over pairs and conditions, scaled as $1.96\,\mathrm{sd}/\sqrt{S}$ --- which is
a different noise model from Eq.~\eqref{eq:var}, the estimator every table in this paper uses. Eq.~\eqref{eq:var}
pools per-configuration run variances and gives per-pair three-seed half-widths of $0.075$--$0.088$
on the same census (Sect.~\ref{subsec7_6}). The two answer different questions --- how far a
whole-census result moves between seed sets, against how precisely one census estimates the benchmark
value --- and we quote each where it belongs rather than treating them as one quantity. A different operating
system, a software-rendered rasteriser, a different graphics path and an out-of-process policy server
move the published numbers by less than one seed's worth of noise, which is the external-validity
check the rest of the paper rests on. On the second embodiment the same check gives $+0.003$
(octo-base on pick-coke-can, 0.173 against 0.170 published) and $-0.040$ (rt-1-x, 0.527 against
0.567).

But the differences are \emph{one-signed}, and that is worth reporting even though we cannot turn it
into a platform constant. Six of the eight cells fall below the published value, two match it, and
none exceeds it; the mean is $-0.038$. We resist three stronger readings that the numbers invite.

The eight cells are four tasks crossed with two policies that share seeds, so they are not eight
independent measurements of one platform, and a $t$ interval over them does not estimate a quantity
that exists. The published rates are themselves point estimates without intervals, so part of the
difference belongs to them. And the lifecycle comparison is suggestive, not causal: our earlier
re-seeding run had a mean of $-0.005$ and a mean across-run range of $0.109$ against $0.052$ now,
and while a noisier estimator agreeing more closely with a target is certainly not evidence of
fidelity, three runs per cell is too few for those ranges to support a claim that extra variance
\emph{concealed} a real gap. What we can say is that the closer agreement we previously reported came
from a protocol we now know to be wrong, and that the corrected protocol agrees less well and in one
direction.

On magnitude, the arithmetic should be stated rather than rounded into a slogan. The mean absolute
difference $0.038$ is about a third of the $0.121$ torque shift, not an order of magnitude below it,
and the largest single cell, $0.111$, is close to it; against the $0.400$ texture range the ratio is
about ten. So the pipeline check licenses less than a sweeping claim of negligibility. What it does
license is the specific use we make of it: every effect in Sects.~\ref{sec5}--\ref{sec8} is a
\emph{difference between two policies measured in the same configuration on our own pipeline and one
build}, where an offset common to both policies cancels and only its interaction with the policy pair
would matter. The build condition is not decoration: Sect.~\ref{subsec5_5} measures a build change
moving $\hat\Delta$ by $-0.109$, so it is not such an offset. The
$-0.038$ is a caution about comparing our absolute rates with published absolute rates, which we do
only in this section.

\subsection{What the published budget resolves}\label{subsec6_3}

The same table answers a second question, and the answer is less comfortable. Even under the
reference lifecycle the three runs disagree with each other by more than we differ from the published
value: the across-run range of a cell averages 0.052 and reaches 0.125. Applying the census estimator
of Sect.~\ref{subsec5_4} to the comparison the published protocol is built to support gives
Table~\ref{tab5}.

\begin{table*}[t]
\caption{What the published 72-episode budget resolves, under the reference RNG lifecycle. Census
interval: Eq.~\eqref{eq:var}, policy noise only, $S = 3$ runs over a complete 24-configuration
enumeration. The last two columns give our $\Delta$ and verdict from the per-episode-re-seeding run
for comparison; the verdict column is unchanged by the lifecycle, the carrot sign is not. Generated
by \texttt{scripts/make\_table5.py}, which reproduces both.}\label{tab5}
\begin{tabular}{@{}lrllrrl@{}}
\toprule
 & \multicolumn{3}{c}{reference lifecycle} & & \multicolumn{2}{c}{re-seeding} \\
\cmidrule(lr){2-4}\cmidrule(l){6-7}
Task & our $\Delta$ & census 95\% & verdict at this budget & published $\Delta$ & $\Delta$ & verdict \\
\midrule
eggplant & $+0.097$ & $[-0.042, +0.236]$ & abstain & $+0.138$ & $+0.083$ & abstain \\
spoon & $+0.361$ & $[+0.259, +0.463]$ & octo-small $\succ$ octo-base & $+0.347$ & $+0.347$ & octo-small $\succ$ \\
carrot & $+0.028$ & $[-0.039, +0.094]$ & abstain & $+0.014$ & $\mathbf{-0.056}$ & abstain \\
stack & $+0.028$ & $[-0.011, +0.066]$ & abstain & $+0.042$ & $+0.042$ & abstain \\
\bottomrule
\end{tabular}
\end{table*}

Three observations.

\paragraph{(i) One of four orderings is resolvable.}
Only the spoon comparison excludes zero. On the other three tasks the published protocol reports a
rate difference its own budget cannot distinguish from noise. This is not a criticism of the
reference implementation --- the published tables report rates, not rankings --- but readers do
compare those rates across policies, and on three of four tasks that comparison is unsupported.

\paragraph{(ii) A sign reversal we previously reported does not survive the corrected lifecycle, and
we withdraw it.}
Our earlier reproduction put octo-base above octo-small on carrot ($-0.056$) where the published
values put octo-small above octo-base ($+0.014$), and we reported that as a concrete instance of a
published ranking failing to reproduce. Under the reference RNG lifecycle carrot gives $+0.028$ --- the
same sign as published --- and no cell in Table~\ref{tab5} disagrees in sign with the published
table. The reversal was an artefact of our own per-episode re-seeding, which gave all 24
configurations in a run one shared noise realisation and so let a single unlucky stream decide the
sign of a $0.014$-sized difference. We therefore no longer claim that a faithful re-run of the
published protocol reverses a published ordering; what we observed was our own implementation
reversing it.

The underlying point survives in the weaker and better-supported form of observation (i): carrot's
difference is $+0.028$ with interval $[-0.039, +0.094]$, so the published $+0.014$ is a number whose
sign its own budget does not determine. A comparison that cannot be resolved can come out either way,
and in our first run it did --- which is a fact about the budget, demonstrated on our own data rather
than charged against the reference implementation.

\paragraph{(iii) Removing configuration sampling error is not enough.}
The census interval is narrower than the binomial-over-72 a reader of the published table would
compute on three of the four tasks and marginally wider on the fourth: half-widths 0.139 against
0.163 on eggplant, 0.102 against 0.115 on spoon, 0.067 against 0.083 on carrot, and 0.039 against
0.038 on stack. No verdict changes either way. With three runs per configuration the
per-configuration policy variance dominates, so enumerating the configuration grid, which sets the
sampling error to exactly zero, buys much less than practitioners expect. Budget is spent better on
seeds than on episodes.

\paragraph{(iv) The verdict count depends on the variance model, and we say which we use and why.}
``One of four'' is not model-free. Eq.~\eqref{eq:var} treats each configuration's runs as the
replication unit and pools the per-configuration policy variances; on the new records that gives the
intervals in Table~\ref{tab5} and one resolvable ordering. An alternative that treats each of the
three runs as the unit --- a paired $t$ interval on three whole-census block differences --- gives
eggplant $[+0.038, +0.157]$, spoon $[+0.242, +0.481]$, carrot $[-0.188, +0.243]$, stack
$[-0.032, +0.088]$, hence \emph{two} resolvable orderings. The block interval is not automatically
better: with three blocks it rests on two degrees of freedom, and eggplant's three block differences
happen to fall close together, which is what makes that interval narrow.

We report Eq.~\eqref{eq:var} as the headline because it matches the randomisation the protocol
actually performs. Under the reference lifecycle each episode draws from its own point in one
advancing stream, so the noise varies within a run and not only between runs, and a model whose
replication unit is the whole run attributes within-run variation to nothing. Under our earlier
re-seeding lifecycle the opposite was true and the block model would have been the right choice ---
which is a further consequence of the lifecycle error, and the reason we state the model explicitly
instead of inheriting it. The sensitivity is reported here rather than buried because the headline
claim is one of the paper's three, and a reader is entitled to know that a defensible alternative
model makes it two of four rather than one.

\subsection{What this section licenses}\label{subsec6_4}

The published numbers carry an uncertainty of roughly $\pm 0.08$ at three seeds, which the published
tables do not quantify, of the
same order as the parameter-induced shift of 0.121 that Sect.~\ref{sec7} attributes to a
calibration-invisible torque limit. That these two are comparable is the paper's central quantitative
claim, and Sect.~\ref{subsec6_2} establishes that we can say it about the benchmark as published, not
merely about our port of it.

Two qualifications travel with that licence. Our rates average $0.038$ below the published
ones, so a comparison of an absolute rate of ours against a published absolute rate should carry that
offset; every claim in Sects.~\ref{sec7} and~\ref{sec8} is instead a \emph{difference} between two
policies measured in the same configuration and on one build, where an offset common to both
policies cancels --- a build change is not one of those (Sect.~\ref{subsec5_5}). And the RNG lifecycle is
part of a protocol, not an implementation detail beneath it: it changed our agreement with the
published table, halved the across-run spread, and decided the sign of one comparison. A reproduction
that matches an outer loop while differing in how the policy's randomness is threaded through it is
not a reproduction of the same estimand, which is a small instance of this paper's general point
applied to our own work.

%%=============================================================================%%
\section{Decision-relevant non-identifiability}\label{sec7}

Sections \ref{sec4} and \ref{sec5} showed that the calibration protocol cannot see certain parameter
directions and that the benchmark's own budget leaves its rankings unresolved. Neither observation is
interesting on its own: an unidentifiable parameter matters only if moving it inside the compatible
set changes an answer someone acts on. This section establishes that it can, measures how large the
effect is, and reports how much of it survives when the evaluation budget is increased --- which is
less than we first claimed.

\subsection{The claim, stated so it can fail}\label{subsec7_1}

Fix a task, a policy pair $(i,j)$ and the invisible fibre $F_{\mathrm{nom}}$ of Sect.~\ref{subsec4_4}. The
claim under test is that point calibration and the union bound disagree --- that point calibration
declares orderings the calibration data cannot support. Two forms are worth separating.

The \textbf{strong form}: some $z \in F_{\mathrm{nom}}$ reverses a ranking's sign, so that the simulator's
answer to which policy is better depends on a parameter the calibration data is silent about.
\textbf{We report this as false.} Across every pair, task and condition we measured, no
calibration-invisible parameter reverses a sign-determined ranking. We also once believed we had
found a weaker version of it --- a between-condition sign change on one pair --- and withdrew that
too (Sect.~\ref{subsec7_5}).

The \textbf{weak form}: the parameter moves $\Delta$ far enough to cross the decision boundary,
turning a declared ordering into an abstention. This holds, and it is budget-dependent --- which is
the finding that replaced our original, stronger claim.

\subsection{One pair, end to end}\label{subsec7_1b}

Before the aggregate counts, one pair with the whole chain in a single place, because the aggregate
tables invite a splice we want to foreclose. This paper reports the largest parameter effect on one
policy pair, a verdict change on a second, and a third arm's published real rates; those are three
measurements and not one closed loop, and Table~\ref{tab6a} is the loop that does exist.

\begin{table*}[t]
\caption{\texttt{octo-small} versus \texttt{octo-base} on the ManiSkill3 eggplant task, followed
from the calibration protocol to the verdict. The policy columns share one scope: all 64 configurations
enumerated, so configuration sampling error is exactly zero; five seed sets; one inference-stack
build; the run observation unit; Eq.~\eqref{eq:var} with empirical variance on both sides.
\textbf{The second and fourth columns are not commensurable and that is the point.} The replay-loss
column is what the calibration protocol can resolve about the condition, in the composite units of
Eq.~\eqref{eq:loss}; \emph{``unobserved''} means the parameter never enters a free-space replay
trajectory, so there is nothing to measure --- a gap in the protocol's coverage rather than in its
precision. The $\hat\Delta$ column is how far the policy comparison moves there, in probability
points. Replay-loss bounds are from the original ManiSkill2 replay sweeps; policy differences and
intervals are from the ManiSkill3 census. The table therefore combines calibration and policy
evidence from different stacks, not a matched single-stack measurement of both columns. In
particular, the sixteenfold gain-scaling replay bound is from the original-stack sweep described in
Sect.~\ref{subsec4_4}; it is not asserted to bound an unmeasured ManiSkill3 sweep over the same range.
Generated by \texttt{scripts/make\_evidence\_chain.py}.}\label{tab6a}
\small
\begin{tabular}{@{}lllrll@{}}
\toprule
Condition & Calib.\ sees it? & Replay-loss change & $\hat\Delta$ & 95\% interval & Verdict \\
\midrule
\texttt{nominal} & base point & $0$ by construction & $+0.055$ & $[-0.008, +0.117]$ & abstain \\
\texttt{iso\_x0.25} & measured flat & $\le 3.1\times10^{-6}$ & $+0.092$ & $[+0.025, +0.159]$ & \texttt{small} $\succ$ \\
\texttt{iso\_x4.0} & measured flat & $\le 3.1\times10^{-6}$ & $+0.174$ & $[+0.105, +0.244]$ & \texttt{small} $\succ$ \\
\texttt{force\_x0.5} & measured flat & $\le 1.8\times10^{-8}$ & $+0.106$ & $[+0.040, +0.172]$ & \texttt{small} $\succ$ \\
\texttt{fric\_x0.4} & unobserved & not in a replay & $+0.077$ & $[+0.015, +0.139]$ & \texttt{small} $\succ$ \\
\texttt{dens\_x0.5} & unobserved & not in a replay & $+0.128$ & $[+0.059, +0.197]$ & \texttt{small} $\succ$ \\
\midrule
\textbf{union over the fibre} & --- & --- & --- & $\mathbf{[-0.008, +0.244]}$ & \textbf{abstain} \\
\bottomrule
\end{tabular}
\end{table*}

Read the two middle columns against each other. Across the fibre $\hat\Delta$ ranges over
$0.120$ --- from $+0.055$ to $+0.174$, which is the difference between an abstention and a
comfortable declaration --- and it does so over conditions whose replay-loss difference from nominal
is at most $3.1\times10^{-6}$ in composite units, or is not measurable at all because the parameter
never appears in the trajectory. \textbf{The calibration protocol's resolution and the policy
comparison's sensitivity are unrelated in magnitude on this pair}, and that disconnection is the
whole mechanism this paper is about, stated on data where every scope parameter is fixed and printed.

Two things this table deliberately does not do. It does not carry the operating-point result, which
needs a second census at the fitted controller setting and appears as Table~\ref{tab10} --- the same
pair, but at three seed sets, so that the two operating points are compared at a matched budget. And
it does not carry the $0.121$ torque shift of Sect.~\ref{subsec7_3}: that is measured on the
cross-family Octo-versus-OpenVLA pairs, where a large-displacement policy saturates contact torque,
and on this Octo-versus-Octo pair the same parameter moves $\hat\Delta$ by $+0.051$. Quoting the
two together is exactly the splice this table exists to prevent.

\subsection{The core table, and what a larger budget did to it}\label{subsec7_2}

The complete per-pair table --- both verdicts for every policy pair on every task, with the
configurations covered, the runs per policy and the variance model for each --- is
\texttt{results/CORE\_TABLE.md} in the release, generated by \texttt{make\_core\_table.py}. It has 18
rows: 17 bridge policy pairs across the two ports, plus the pick-coke-can pair of
Sect.~\ref{subsec8_4}, which we count separately because it is a different embodiment and a
different question. Wherever this paper says ``of 17 pairs'' it means the 17 bridge rows.
Table~\ref{tab6} below is not that table; it shows only the three pairs whose point-versus-set
verdict differs at the reference protocol's three seed sets, and what raising the budget to five
does to each.

Before the numbers, the observation unit, because the estimator has two defensible ones and they
disagree on two rows. Eq.~\eqref{eq:var} writes $S_c$ for the number of \emph{runs} at
configuration $c$, and that is what we report: episode outcomes pooled across seed-set directories,
so a configuration that received six episodes contributes six observations. An earlier
implementation instead averaged a directory's episodes at a configuration and treated each
\emph{directory mean} as the observation, making $S_c$ the number of directories. The two differ
because seed set B carries 96 episodes on a 64-configuration grid, so ids 64--95 wrap onto
configurations 0--31 and give those two episodes --- at different policy seeds, since
$\text{policy\_seed} = \text{base} + \text{episode\_id}$. Neither unit is wrong: the run unit is
finer and is what the formula states; the directory unit is conservative against a directory-level
common effect, which the build-drift result of Sect.~\ref{subsec5_5} shows can exist across
generations, though the five sets here share one build. We adopt the run unit because it is the
formula's, and report the other as a sensitivity --- not the reverse, and not chosen by which
declares more. \texttt{make\_core\_table.py --unit directory} reproduces it. The two units move two
of the eighteen rows: \texttt{octo-small} vs \texttt{octo-base} on eggplant declares under the
directory unit on a lower bound of $+0.000042$ and abstains under the run unit at
$[-0.008, +0.117]$, and \texttt{octo-small} vs \texttt{octo-base} on the original-stack spoon does
the opposite. \textbf{The point-versus-set disagreement count is 1 of 17 under both.}

Of the 17 bridge pairs, \textbf{one} disagrees: \texttt{octo-base} vs \texttt{octo-base@hist1} on the
eggplant census, point $\Delta = +0.073\,[+0.01, +0.14]$ declaring \texttt{octo-base} better, union
interval $[-0.06, +0.14]$ abstaining.

The union runs over the nominal setting as well as the five calibration-invisible ones, and the
reason is definitional rather than statistical. The set being unioned over is $F_{\mathrm{nom}}$,
the fibre of settings the replay loss cannot distinguish from nominal (Sect.~\ref{subsec4_4}), and
nominal is its base point, so it belongs to the fibre by construction. An envelope taken over the
perturbed conditions alone would not cover the fibre it claims to bound. It widens no verdict in
Table~\ref{tab6}. We note that nominal's membership in $F_{\mathrm{nom}}$ is a tautology and carries
no implication that nominal is the calibration optimum --- Sect.~\ref{subsec4_4} shows it is not.

One consequence should be stated plainly, because a reader scanning the table's two-decimal columns
cannot see it. Two declarations rest on a bound that rounds to $+0.01$: \texttt{octo-small} vs
\texttt{octo-small@hist1} at $+0.0091$, and the surviving disagreement's own point calibration at
$+0.0095$. Put a lower bound of $+0.0095$ in units of the evidence behind it. On a 64-configuration
census at five runs per policy, $\hat\Delta$ weights each configuration mean by $1/64$ and each
\emph{individual episode} by $1/(64 \times 5) = 0.0031$, so $+0.0095$ is about $0.6$ of one
configuration's mean outcome, or about \textbf{three single-policy episodes} out of the 320 on each
side of the comparison. Three episodes flipping the other way would erase it.
Neither is a declaration in any useful sense; both sit on the knife edge, and we flag them rather
than let the rounding carry them. The criterion's substantive declarations are the ones with margin,
and the multiplicity analysis below says which those are: the five cross-family Octo-versus-OpenVLA
envelopes that survive a family-wise correction, plus the fractal pair of Sect.~\ref{subsec8_4}. The
spoon envelopes are not among them --- $[+0.04, +0.56]$ looks wide but its \emph{lower} bound is
$+0.035$, and the envelope's binding constraint is that bound rather than the width of the interval,
so it does not survive. That is worth noticing: a wide envelope that only just clears zero is a weak
declaration, not a strong one, and reading the interval's width for its margin gets it backwards. The
generator marks every knife-edge row, so the distinction cannot be lost in a later revision. We also
report, rather than bury, that the knife edges are not stable across the two observation units ---
the largest of them under the directory unit was $+0.000042$, on a pair that abstains under the run
unit. A bound that small is a property of the estimator's last digit, not of the policies.

At three seed sets, three rows disagreed. The two that no longer do were resolved by spending seeds,
not by any change of method or of data-retention rule.

\paragraph{The counts above are uncorrected for multiplicity, and under a common error-control
standard the two rules disagree \emph{more}, not less.}
Seventeen pairs are seventeen hypotheses, and a reader who scans the table for whichever orderings
it declares is making a family-wise claim. We therefore apply Holm's step-down procedure at
$\alpha = 0.05$ to each rule's 17 $p$ values --- Holm needs no independence assumption, and these
pairs are certainly not independent, six of them being drawn from four Octo variants on shared
policy data (\texttt{scripts/analyze\_multiplicity.py}).

\begin{table}[h]
\caption{The same 17 bridge pairs under two error-control standards. The last column is the
quantity this section reports.}\label{tab6b}
\begin{tabular}{@{}lrrr@{}}
\toprule
Standard over the 17 pairs & point & envelope & disagree \\
\midrule
no correction across pairs & 11 & 10 & \textbf{1} \\
Holm at $\alpha = 0.05$ & 9 & 5 & \textbf{4} \\
Benjamini--Hochberg & 11 & 10 & \textbf{1} \\
\bottomrule
\end{tabular}
\end{table}

\noindent Read the last column, because it is the one the paper's claim is about. Point
calibration's 11 declarations fall to 9 --- the two that go are exactly the knife edges of the
previous paragraph, at $+0.0091$ and $+0.0095$ --- and the envelope's 10 fall to 5, so the
\emph{disagreement} between the rules rises from 1 of 17 to \textbf{4 of 17}. Quoting the envelope's
own fall from 10 to 5 in isolation, as an earlier draft of this paragraph did, drops the comparison
the section exists to make: correcting for multiplicity makes the set-valued rule \emph{more}
consequential relative to point calibration, not less, because it removes more of the envelope's
declarations and so turns agreements into disagreements. The four pairs are
\texttt{octo-small} vs \texttt{octo-base@hist1} and \texttt{octo-base@hist1} vs \texttt{openvla} on
the eggplant census, and \texttt{octo-small} vs \texttt{octo-base} on spoon on each stack.

Three limits on that reading, and we hold to all three. \textbf{Four extra abstentions are not four
corrected errors.} We have no true ordering for these pairs, so we cannot tell an avoided false
declaration from a surrendered correct one, and the two rules are not even answering the same
question --- one asks about a single setting, the other about agreement across a set. What is
established is that a declaration is sensitive to the parameter set and to the error-control
standard; that a set-valued rule \emph{decides better} is not established here and we do not claim
it. \textbf{Holm needs valid input $p$ values}, which ours are only under the normal approximation
and the variance imputation of Sect.~\ref{subsec5_4}; a correction does not repair a calibration
problem in the thing it corrects, so this is a sensitivity analysis under a stated working model
rather than a guarantee. And \textbf{the direction of the envelope's larger loss is a fact about
these records, not a theorem}: an intersection--union test's $p$ value is the largest of its
components, which explains why the envelope declares less often to begin with, but it does not
determine how a step-down procedure over a different index --- pairs, not conditions --- will treat
it. None of this touches the no-correction-across-conditions argument of Sect.~\ref{subsec3_3}, and
a correction cannot affect an abstention, since it only removes declarations. Benjamini--Hochberg
leaves every count unchanged; we report Holm because the reading at issue is family-wise rather than
one about a false-discovery rate.

\begin{table*}[t]
\caption{The effect of raising the seed budget from three to five on the point-versus-set
disagreements of the eggplant census. Both columns are recomputed on the run observation unit with
$S=3$ defined as seed sets A$'$, B and C, so the two are directly comparable and both are
reproducible from the release with \texttt{make\_core\_table.py -}\texttt{-max-sets 3}. The third
row's $S=3$ abstention rests on a lower bound of $-6\times10^{-6}$, which is a disagreement in
arithmetic and not in substance; we list it because the count is ours to report, not to curate. Envelopes include the nominal setting, as the bound requires.}\label{tab6}
\begin{tabular}{@{}llll@{}}
\toprule
Pair & Union at $S=3$ & at $S=5$ & Outcome \\
\midrule
octo-small vs octo-base@hist1 & $[-0.020, +0.275]$ abstain & $[+0.012, +0.245]$ declare & withdrawn \\
octo-base@hist1 vs openvla & $[-0.000006, +0.416]$ abstain & $[+0.017, +0.426]$ declare & withdrawn \\
octo-base vs octo-base@hist1 & $[-0.099, +0.167]$ abstain & $[-0.064, +0.137]$ abstain & survives \\
\bottomrule
\end{tabular}
\end{table*}

We had pre-committed to this test and to withdrawing a row if spending seeds resolved it, which is
why the surviving count is one and not three.

The honest general statement is therefore about the evaluation budget. At the reference protocol's
three seed sets, point calibration declares orderings the union bound does not support on
\textbf{3 of the 17} bridge policy pairs --- though one of the three, as the caption says, abstains
by six millionths. At five seed sets over a complete configuration census, one remains. A single
seed set can evaluate only 13 of the 17 pairs at all, and disagrees on 2 of those 13. The
criterion's value is concentrated in the budget
regime practitioners are actually in. This is weaker than the claim we set out to make, and it
remains useful: we are aware of no published evaluation that runs five seed sets over a full
configuration census, and the reference protocol itself uses three. Every count here is reproducible
with \texttt{make\_core\_table.py -}\texttt{-max-sets}, which is the option we added when we found
that the historical counts could not be recomputed.

The surviving pair is also instructive about \emph{why} it survives. It is not that its $\Delta$
changes sign across conditions (Sect.~\ref{subsec7_5}), but that its nominal interval is barely
decidable (lower bound $+0.0095$) while the union bound takes the minimum lower bound over
conditions, which friction $\times 0.4$ supplies at $-0.064$. It is a margin-and-width phenomenon,
not a parameter-sign phenomenon --- and one whose point declaration sits on the knife edge, so the
disagreement is between an abstention and a barely-made declaration rather than between two
confident answers.

\subsection{The mechanism: one parameter, measured on a complete census}\label{subsec7_3}

The torque limit is the cleanest case. It is invisible to calibration
(Sect.~\ref{subsec4_2}(ii)) because along demonstrated velocities the commanded torque does not
saturate, so replay loss does not depend on it: bitwise on 97 of the 98 demonstrations, and to
0.050~mm on the one that evidently comes closest to the limit. Halving it therefore stays inside
$F_{\mathrm{nom}}$ with a margin set by the mechanism rather than by a fitting tolerance. On the 64-configuration eggplant census it
moves the cross-family comparison by $-0.104$, $-0.146$, $-0.108$ and $-0.127$ for the four
Octo-versus-OpenVLA pairs: a shift of 0.104--0.146, mean \textbf{0.121}, uniform enough across pairs
to be consistent with a single policy moving rather than four independent effects
(\texttt{scripts/analyze\_torque\_shift\_s5.py}).

\begin{remark}
Two earlier measurements of this shift were larger, and both were superseded for the same kind of
reason. On 48 of the 64 configurations it read 0.23--0.25, with two pairs crossing the decision
boundary: the partial census overstated it. On the complete census at three seed sets it read
0.133--0.161, mean 0.141; at five it reads as above. Every effect size in this paper is census-scope
at $S = 5$ for this reason, which is the budget Sect.~\ref{subsec7_2} establishes as the one that
settles verdicts.
\end{remark}

\subsection{Why it is not an artefact}\label{subsec7_4}

Three checks.

\paragraph{It is a property of one policy, and a large one.}
Halving the torque limit raises OpenVLA's eggplant benchmark value from 0.156 to 0.271 on the
complete 64-configuration census ($+0.115$), while octo-small, run as a negative control over the
same census, moves by $+0.011$ and the other three Octo policies by $-0.06$ to $+0.02$ with no
consistent direction. The asymmetry has a mechanism: OpenVLA emits larger single-step displacements,
so at the nominal limit it pushes the eggplant away or re-grasps repeatedly on contact, and the
episode-level counters agree (grasp rate $+0.12$/$+0.18$, fewer repeated gripper closures).

\paragraph{It is broad, not driven by a few configurations.}
Over the complete census 16 configurations improve, 42 are unchanged and 6 degrade, and the mean
change is positive in 7 of the 8 eggplant orientations ($-0.04$ to $+0.29$). The gain is
anticorrelated with nominal success ($r = -0.44$): of the 51 configurations where OpenVLA fails
outright at the nominal limit, 13 are rescued. The parameter repairs failures rather than padding
successes. The same decomposition for octo-small is flat and directionless --- 15 improve, 15
degrade, 2 of 8 orientations positive --- which is what a policy that never saturates the limit
should look like.

\paragraph{It is task-dependent in the direction the mechanism predicts.}
On the spoon task (24-configuration complete census) the same parameter barely moves $\Delta$ at all
--- $-0.01$ to $-0.07$ across the four cross-family pairs, against $-0.14$ on eggplant --- and
changes no verdict. Eggplant is where contact torque saturates for a large-displacement policy;
spoon is not. A calibration-invisible parameter is therefore not uniformly dangerous: its
decision-relevance is a property of the task-policy pair, which is precisely why it cannot be argued
away a priori and has to be enumerated.

\begin{figure*}[t]
\centering
\includegraphics[width=\textwidth]{Fig3}
\caption{The torque-limit effect decomposed over the eggplant orientation grid, for a
deterministically decoding policy (OpenVLA, left) and a stochastic one (Octo-small, right). The
effect is positive in 7 of 8 orientations for the former and directionless for the latter, which is
the asymmetry the saturation mechanism predicts.}\label{fig3}
\end{figure*}

\subsection{A withdrawn claim: there is no between-condition sign reversal}\label{subsec7_5}

An earlier version of this section reported that one pair, \texttt{octo-base} vs
\texttt{octo-base@hist1}, reversed sign across the invisible fibre --- $\Delta = -0.021$ under
density $\times 0.5$ against positive values under the other five conditions --- and argued from it
that this pair's abstention was a fact about the parameters rather than about our budget, and so
could never be removed by more evaluation. \textbf{We withdraw that.} At five seed sets the same
cell reads $+0.042\,[-0.022, +0.106]$, and no condition has a point estimate below $-0.001$:
nominal $+0.069$, force $\times 0.5$ $+0.050$, iso $\times 0.25$ $+0.030$, iso $\times 4.0$
$+0.006$, friction $\times 0.4$ $0.000$, density $\times 0.5$ $+0.042$. The between-condition sign
reversal we had reported is not merely smaller at five seed sets; it is absent.

The distinction the withdrawn claim was reaching for is nevertheless real and worth keeping, because
Sect.~\ref{sec8} needs it. As the evaluation half-width $h \to 0$ the union bound converges to
$\big[\min_c \Delta_c, \max_c \Delta_c\big]$, which does \textbf{not} shrink with budget. An
abstention that survives $h \to 0$ is a parameter fact; one that dissolves as $h$ falls was a power
artefact. Our own three-to-five-seed experiment is the clean demonstration of the second case, and
Sect.~\ref{subsec8_4} uses the same test to reject an abstention we would otherwise have been
tempted to claim as a success.

\subsection{What determines whether a ranking can be declared}\label{subsec7_6}

A ranking can be declared when $|\hat\Delta|$ exceeds the parameter-induced shift $\delta$
\textbf{plus} the evaluation half-width $h$. Enumerating the configuration grid drives the sampling
part of $h$ to zero; more seeds shrink the rest, as $1/\sqrt{S}$ --- from 0.075--0.088 at three seed
sets to 0.058--0.068 at five. The parameter-induced shift $\delta \approx 0.121$ does not move, by
construction.

The two regimes this produces are the practical content of the section. At the budget the reference
protocol uses --- 24 configurations, three seeds --- a pair's half-width is 0.05--0.13
(Sect.~\ref{subsec6_3}), which is at or above $\delta$; there the two error sources are not even
separable, and a practitioner cannot tell an identifiability problem from a power problem. Our census
at $S = 5$ pushes $h$ below $\delta$, and only then does the comparison become meaningful. The
recommendation follows directly: report the estimand and the budget, and spend budget on seeds
before episodes, because only then does the residual ambiguity that \emph{cannot} be bought down
become visible.

\subsection{The ranking at the calibration-preferred operating point}\label{subsec7_7}

Everything above is computed at the simulator's shipped nominal controller setting, which
Sect.~\ref{subsec4_4} showed our re-application of the calibration protocol disfavours in favour of
$(k{\times}2,\ d{\times}0.5,\ \text{delay}\ 1)$. That leaves an obvious question, and rather than
argue about it we measured it: we re-ran the complete 64-configuration eggplant census for
\texttt{octo-small} and \texttt{octo-base} at the fitted setting, under three seed sets whose bases
match three of the census sets used throughout. The comparison is paired in everything that can be
held fixed --- the same 64 configurations, the same policy seeds episode for episode, the same
policies, the same build, the same estimator (Eq.~\eqref{eq:var}) --- so the only difference between
the two columns of Table~\ref{tab10} is where the controller was set. Nominal is restricted to the
same three seed sets, because comparing against the $S=5$ interval would confound the operating
point with the budget.

\begin{table*}[t]
\caption{The same policy pair, the same 64 configurations and the same policy seeds, evaluated at
the benchmark's operating point and at the one its replay objective prefers. \textbf{The point
verdict changes and the set verdict does not.} The replay-loss column is the composite objective of
Sect.~\ref{subsec4_1}, in its own units, not a distance --- see Sect.~\ref{subsec4_4}. The set
envelope at each operating point is over that point's own complete six-condition fibre --- six
64-configuration censuses at each setting, 12 in total for this table, all complete. Generated by
\texttt{scripts/analyze\_fitted\_point.py}, which refuses any condition that is not a complete
census.}\label{tab10}
\begin{tabular}{@{}lrlll@{}}
\toprule
 & Replay loss & & \multicolumn{2}{c}{verdict} \\
\cmidrule(l){4-5}
Operating point & vs minimiser & $\hat\Delta$ & point $95\%$ & set envelope \\
\midrule
nominal, $d/k = 1$, no delay & $+0.0029$ & $+0.073$ & $[-0.016, +0.161]$ & $[-0.035, +0.300]$ \\
 & (rejected) & & abstain & abstain (6 cond.) \\[2pt]
fitted, $d/k = 0.25$, delay 1 & $0$ & $+0.104$ & $[+0.016, +0.192]$ & $[-0.029, +0.219]$ \\
 & (the minimiser) & & \texttt{small} $\succ$ \texttt{base} & abstain (6 cond.) \\
\bottomrule
\end{tabular}
\end{table*}

\paragraph{The verdict changes; the underlying difference is not resolved.}
This distinction is the whole content of the result and we lead with it. What we observe is a
\emph{verdict} change on this batch: at nominal the interval includes zero and the criterion
abstains, at the fitted point it excludes zero and declares. What we do \emph{not} observe is a
resolved change in the policy difference. Both $\hat\Delta$ are positive --- there is no sign
reversal --- and the paired change itself,
$D = \hat\Delta_{\mathrm{fitted}} - \hat\Delta_{\mathrm{nominal}} = +0.031$, is not distinguishable
from zero under any sampling model we can defend.

The primary one has to match this paper's own estimand, and that settles it: the 64 configurations
are enumerated and fixed, so they are not a sampling unit, and what remains random is the policy
seed. Each seed set supplies all four arms --- both policies at both operating points --- at every
configuration, so it yields one complete paired shift
$d_c^{(k)} = [\hat p_A^{\mathrm{fit}} - \hat p_B^{\mathrm{fit}}] - [\hat p_A^{\mathrm{nom}} - \hat p_B^{\mathrm{nom}}]$
there, and $\operatorname{Var}(D) = N^{-2}\sum_c \operatorname{Var}_k(d_c^{(k)})/K$ keeps the pairing
rather than assuming the arms independent. That gives
$D \in \mathbf{[-0.086, +0.148]}$. Three sensitivities bracket it: summing the four arms' run
variances separately, which discards the pairing, gives $[-0.093, +0.156]$; treating configurations
as a sampling unit --- a superpopulation reading, and the estimand
Sect.~\ref{sec5} argues is not the benchmark's --- gives $[-0.064, +0.126]$; whole seed sets as
blocks on two degrees of freedom
give $[-0.300, +0.363]$. The three seed sets individually give $+0.172$, $+0.016$ and
$\mathbf{-0.094}$ --- one of the three moves the other way. All four intervals contain zero and none
is an exact finite-sample guarantee. So a significant interval at one operating point and a
non-significant one at the other does not make the difference between them significant, and we do
not claim a stable population effect.

What the measurement does establish is fragility, and that is the paper's thesis rather than a
weaker version of it. A change of $0.031$ --- less than half the $\pm 0.088$ half-width at this
budget, a quarter of the $0.121$ torque shift, and itself unresolved --- moves the answer from
abstention to a declaration, because at nominal the lower bound already sat at $-0.016$. The
declaration at the fitted point rests on a lower bound of $+0.016$. This comparison is matched at
three seed sets, so each policy contributes $64 \times 3 = 192$ episodes and each weighs $0.0052$ in
$\hat\Delta$: the bound is about \textbf{three single-policy episodes}, or one configuration's mean
outcome out of 64. Three episodes are what separates the declaration from an abstention. So at the budgets in use the
verdict here turns on the operating point --- and, unlike the modelling choices of
Sect.~\ref{subsec7_2}, this is a choice the calibration evidence \emph{does} settle, against the
setting the benchmark uses. That is the C3(ii) case rather than the C3(i) case, and we keep them
apart: what is unsettled is not which point the evidence prefers but which verdict is correct, since
the difference between the two points is not itself resolved. A reader given only the published rates
cannot see any of this, because those rates exist at one operating point and carry no interval.

\paragraph{The set-valued verdict does not change: it abstains at both operating points.}
Taking the union over each point's own calibration-invisible fibre gives $[-0.035, +0.300]$ at
nominal and $[-0.029, +0.219]$ at the fitted point, each over that point's complete six-condition
fibre --- twelve 64-configuration censuses in all. Both abstain.

One methodological note, because it cost us a wrong claim and the mistake is easy to make. An
earlier version of this section reported the fitted envelope over the first two completed
conditions, where it was $[+0.009, +0.192]$ and declared, and concluded that both verdicts flip. The
third condition --- a common rescaling of the fitted gains --- gives $\hat\Delta = +0.057$ with
$[-0.029, +0.144]$, which pulls the envelope across zero on its own. An envelope reported before its
design is complete is not a conservative version of the final answer but a different and narrower
one, and here it pointed the opposite way. The monotonicity that makes the completed envelope safe
to report --- a min/max over conditions only widens --- is exactly what makes a partial one unsafe.

Both fibre memberships are measured rather than assumed, which took two additional replay
experiments. The common-scale directions: the objective is flat along them in 14 of the replay
grid's 26 $(\text{ratio}, \text{delay})$ groups, across seven ratios from $0.5$ to $2.0$, to within
$7.2\times10^{-6}$ in composite units, but ratio $0.25$ has exactly one grid point per delay value,
so we swept the common scale sixteenfold at ratio $0.25$ directly and found it flat to
$5.4\times10^{-6}$ on ManiSkill3 and $1.1\times10^{-5}$ on the original stack. The torque limit: the
argument at nominal is that the commanded torque never reaches the limit, and that does \emph{not}
transfer for free, because commanded torque depends on the gains and the fitted point doubles the
stiffness. Measured at the fitted point, halving the limit leaves the replayed trajectory bitwise
unchanged on 96 of 98 demonstrations against 97 of 98 at nominal, and on the demonstrations that do
change the largest end-effector position difference is $84.5$~\textmu m against $50.1$~\textmu m ---
so the clip does activate a little more often and a little harder, as the mechanism predicts, while
remaining far below the protocol's own resolution. Both directions are retained at the fitted point
on evidence rather than by inheritance. What the records cannot show is the mechanism itself:
commanded torques were not logged, so ``the clip activates more often'' remains the natural reading
of these counts rather than something we measured directly.

One further observation, which corrects something we had stated too strongly. The union bound's
abstention at both operating points should not be read as the criterion \emph{protecting} against
the operating-point choice. An envelope over a set cannot speak about a point the set excludes, and
nominal's fibre excludes the fitted point by construction --- the two differ in the ratio and the
delay, which the calibration data constrain, so neither invisible fibre contains the other's base
point. A set containing both operating points and the region between them \emph{would} have an
envelope covering both; it would simply be much wider. So there is no third irreducible kind of
ambiguity here, and we withdraw that framing: it is a question of set definition, and it identifies
a choice the paper had left implicit. The invisible fibre is the right object for asking what the
published numbers determine, and the wrong object for asking whether the operating point itself was
well chosen. What addresses the latter is not a wider envelope but reporting the operating point and
its standing against the calibration objective --- which costs nothing, and which no benchmark we
know of does.

\begin{figure*}[t]
\centering
\includegraphics[width=\textwidth]{Fig4}
\caption{$\Delta$ between two Octo variants by simulator condition, on three bridge tasks. Grey
points are conditions the calibration data identifies, coloured points are the
calibration-invisible ones, and the shaded band is the union bound over the invisible fibre --- the
invisible conditions \emph{and} nominal, which the set contains by construction.
Computed with the same estimator and the same seed sets as Table~\ref{tab6}, so figure and table
report one computation rather than two, on the run observation unit of Sect.~\ref{subsec7_2}.
Verdicts and set bounds: carrot, point and set both abstain, $[-0.20, +0.07]$; spoon, both declare
octo-small better, $[+0.04, +0.56]$; eggplant, both abstain, $[-0.008, +0.244]$. Available seed sets differ
by task --- five for eggplant, two for spoon and carrot --- and the per-cell configuration and run
counts are in Appendix~\ref{secC}.}\label{fig4}
\end{figure*}

%%=============================================================================%%
\section{A verdict ladder and its reach}\label{sec8}

Sections \ref{sec4} through \ref{sec7} establish that a simulation benchmark's ranking claims are
underdetermined by the evidence behind them, from three independent directions. This section asks
what to do about it: what a set-valued criterion buys, at what cost in power, and where it stops
working. We answer all three, and the third answer is negative.

\subsection{Two rungs, with different guarantees}\label{subsec8_1}

\textbf{The union bound} declares an ordering when every \emph{sampled} parameter setting in
$F_{\mathrm{nom}}$ agrees in sign. Its guarantee is of \textbf{sampling-extremum} type: valid with respect
to the extremes of $\Delta(z)$ over the settings actually run, and silent about values between them.

\textbf{A Gaussian-process simultaneous band} fits $\Delta(z)$ over the parameter axes and requires
the sign to hold across a simultaneous band over the whole compatible region. We construct it by
drawing $4{,}000$ posterior functions on a grid over the region, marginalising over a grid of
$(\ell, \sigma_f)$ hyperparameters with the marginal-likelihood weights, and taking $L$ as the
$2.5$th percentile of each draw's \emph{minimum} and $U$ as the $97.5$th percentile of each draw's
\emph{maximum}. Both endpoints are therefore statements about the extremes of a whole function rather
than about $\Delta$ at any one point, which is what makes the band simultaneous over the region and
what makes it expensive in power. Its guarantee is of
\textbf{within-band} type under a smoothness prior, and it therefore covers features that fall
between sampled settings --- at a cost in power that Sect.~\ref{subsec8_2} quantifies. Marginalising
over the hyperparameters rather than fixing them at the maximum-likelihood values matters here: the
13 design points do not pin $\ell$ down, and a band conditioned on one $\ell$ understates the
uncertainty that comes from not knowing it.

An earlier design had a third element, a flatness gate that chose between the two per axis. We
removed it; Sect.~\ref{subsec8_2} explains why.

\subsection{Coverage and power on synthetic ground truth}\label{subsec8_2}

Because real data provide no ground truth about $\Delta(z)$, we built a synthetic track in which the
response surface is known and the criteria can be scored directly: 300 repetitions per cell over
four surface families, crossed with true margins $\Delta_0$ and episode budgets $n$.

\begin{table*}[t]
\caption{Minimum coverage / maximum error rate / mean declaration rate on cells where the truth is
decidable, by response-surface family. The union bound fails exactly where the decisive feature
falls between sampled settings. All three columns are computed on one configuration --- the strip
extended past the sampled design, matching Sect.~\ref{sec4}'s finding that the invariant direction
is unconstrained at every magnitude --- because the strip extent enters the ground truth $[L^*, U^*]$
and therefore every method's coverage, not only the band's. The narrowed-strip variant is reported in
the text rather than mixed into this table (\texttt{scripts/make\_table7.py} prints
both).}\label{tab7}
\begin{tabular}{@{}llll@{}}
\toprule
Surface & Point calibration & Union bound & GP simultaneous band \\
\midrule
Flat & 0.92 / 0.05 / 0.53 & 1.00 / 0.00 / 0.24 & 1.00 / 0.00 / 0.14 \\
Linear & 0.24 / 0.36 / 0.71 & 0.99 / 0.00 / 0.29 & 0.95 / 0.00 / 0.15 \\
Dip at a sampled setting & 0.00 / 0.88 / 0.45 & 0.98 / 0.00 / 0.03 & 0.97 / 0.00 / 0.04 \\
Dip between settings & 0.00 / 0.85 / 0.43 & \textbf{0.22 / 0.17 / 0.21} & 0.91 / 0.00 / 0.08 \\
\bottomrule
\end{tabular}
\end{table*}

\paragraph{(i) Point calibration's error rate rises with the episode budget.}
On any non-flat surface the nominal interval covers $\Delta(\hat z)$ rather than the extremes over
$C_\alpha$, so more episodes shrink the interval around a value that was never the quantity of
interest. On the linear surface at $\Delta_0 = 0.1$ the error rate goes from 0.13 at $n = 24$ to
0.36 at $n = 96$; on the sampled-dip surface at $\Delta_0 = 0.2$, from 0.26 to 0.88. This is the
quantitative form of ``more evaluation makes a wrong conclusion more confident''.

\paragraph{(ii) The union bound attains its guarantee, and only its guarantee.}
Where the extremum of $\Delta(z)$ falls on a sampled setting it has coverage $\ge 0.98$ with error
rate 0 and the highest power of any valid rung (declaration rate 0.98 on the flat surface at
$\Delta_0 = 0.3$, $n = 96$). Where the decisive feature falls between sampled settings it fails
outright: coverage 0.22, error rate 0.17.

\paragraph{(iii) The GP band's \emph{observed} coverage is at least 0.91 on all four families, with
error rate 0, and it pays for that in power.}
We are careful with the word. These are empirical coverage frequencies over 300 repetitions per
cell, not a proved guarantee, and 0.91 is below the nominal 0.95 the construction aims at. The Monte
Carlo standard error at 300 repetitions is about 0.017 near a coverage of 0.9, so the 0.91 cell sits
roughly two standard errors below nominal and the shortfall is unlikely to be simulation noise
alone. The band's validity is \emph{conditional on its smoothness prior}; what Table~\ref{tab7}
reports is how often that prior was good enough on these four surface families, which is evidence
about the families we chose rather than a distribution-free statement.
Its declaration rate on the flat surface at $\Delta_0 = 0.3$, $n = 96$ is 0.47, and barely grows
from $n = 48$ to $n = 96$, because the band is evaluated half a grid step beyond the sampled range
and the posterior diverges at the ends.

Narrowing the strip to the sampled range raises that rate to 0.77 and restores monotonicity in $n$,
and it is tempting to adopt it on that basis. We report it as a sensitivity analysis instead,
because the change is not confined to the band. The strip extent enters the ground truth
$[L^*, U^*]$, so narrowing it narrows what every method is asked to cover: point calibration's
maximum error rate on the linear surface falls from 0.36 to 0.04 and the union bound's mean
declaration rate from 0.29 to 0.19, purely from the redefinition. Comparing the band's 0.91 observed
coverage under one truth against its 0.85 under another is therefore not a like-for-like measurement
of one estimator, and we do not present it as one. Within the narrowed configuration, taken whole,
the band's minimum observed coverage is \textbf{0.85} on the between-settings dip with a nonzero
error rate (0.01) --- further below nominal than the headline table's figure, and we attach the
$\ge 0.91$ observation only to the configuration Table~\ref{tab7} reports.

On the real data of Fig.~\ref{fig6} the situation is different and the choice is forced: the 13
design points reach $|u|, |v| = 2$, and a verdict read off the band beyond that would be
extrapolation with no observation behind it. Fig.~\ref{fig6} therefore evaluates the band over the
sampled range only, and says so.

The highest-weight length scale is $\ell = 0.5$ on all three tasks, with signal scale 0.02 on carrot
and 0.05 on spoon and eggplant, so the band there does not rest on a per-task tuning choice.

One scope limitation has to be stated before Fig.~\ref{fig6} is read alongside anything else in this
paper, because it is easy to miss and it invalidates the obvious comparison. The controller plane
needs 13 design points spanning stiffness and damping independently, and the only sweep we have with
that design is the \emph{pre-fix} one: an older inference-stack build, 48 episodes per condition, and
per-episode rather than per-configuration standard errors. Every other number in this paper is
census-scope on the current build at $S = 5$. Fig.~\ref{fig6} is therefore an \textbf{exploratory
analysis on a different and superseded data scope}, not a rung evaluated on the same footing as
Table~\ref{tab6}, and we do not cross-validate the two against each other. Concretely: on spoon the
GP band over that older sweep declares ($[+0.03, +0.43]$) while the union bound over the \emph{same}
older sweep abstains with a lower bound of exactly $0.000$; on the current census
(Table~\ref{tab6}) the spoon union bound is $[+0.04, +0.56]$ and declares. Those three numbers come
from two different data scopes and one cannot be read as a methodological disagreement with another.
What Fig.~\ref{fig6} contributes is a picture of the response surface over a plane no current sweep
covers --- including the ratio direction that Sect.~\ref{subsec4_4} shows the calibration data care
about --- and that is the only weight we place on it.

We also report a configuration we \textbf{rejected}, because it is a methodological point rather
than a tuning detail. Imposing a lower bound on the GP length scale equal to the sampling spacing
doubles power again (flat-surface declaration 0.22 to 0.40, linear 0.17 to 0.30) but assumes the
surface is smooth at exactly the scale where the dip families are not, and coverage collapses to
0.45--0.83 with error rates up to 0.10. A response-surface method's power can be bought only with
prior assumptions, and the stronger the prior the more brittle the method is to the features it did
not sample. The flatness gate we removed is the same lesson: with 13 design points the GP already
fits the sampled-dip surface, so the gate only lowered coverage on the between-settings surface
(0.91--0.98 to 0.89) and mis-fired on flat surfaces 16--20\% of the time.

\begin{figure*}[t]
\centering
\includegraphics[width=\textwidth]{Fig5}
\caption{The three criteria on synthetic ground truth, by response-surface family (columns), true
margin $\Delta_0$ (horizontal axis) and episode budget (line style). \textbf{Top row}: false
declaration rate, with the nominal 0.05 level marked. \textbf{Bottom row}: coverage of
$[L^*, U^*]$, with the nominal 0.95 level marked. Point calibration (grey) gets \emph{worse} as the
budget grows on every non-flat surface --- its false declaration rate reaches 0.88 at $n = 96$ on the
sampled-dip surface, and its coverage falls towards zero. The union bound (blue) holds at
$\ge 0.98$ coverage except on the between-settings dip, where it drops to 0.22 and declares falsely
at up to 0.17; the GP simultaneous band (orange) attains an observed coverage of at least 0.91 on all
four families, which is below the nominal 0.95 and is a measured frequency, not a guarantee. Drawn from the same cached cells as Table~\ref{tab7} and therefore on the same strip extent;
the narrowed-strip variant discussed in Sect.~\ref{subsec8_2} is not plotted, and neither is the
withdrawn flatness gate.}\label{fig5}
\end{figure*}

\begin{figure*}[t]
\centering
\includegraphics[width=\textwidth]{Fig6}
\caption{Posterior mean of the policy difference over the controller plane, for three bridge tasks,
marginalised over the GP hyperparameter grid; circles are the 13 sampled design points and the
dashed lines bound the replay-compatible strip (gain ratio $\approx 1$). The surface drawn and the
bounds quoted are one computation --- the restricted simultaneous band adopted in
Sect.~\ref{subsec8_2}, evaluated over the sampled range $|u|, |v| \le 2$ with 4{,}000 posterior
function draws. Bounds and verdicts: carrot $[-0.10, +0.03]$, abstain; spoon $[+0.03, +0.43]$,
octo-small better; eggplant $[-0.06, +0.36]$, abstain. The highest-weight
hyperparameters are $\ell = 0.5$ with $\sigma_f = 0.02$ (carrot) and $0.05$ (spoon, eggplant).
\textbf{Data scope: this figure uses the older pre-fix sweep} (48 episodes per condition, per-episode
standard errors), the only design that varies stiffness and damping independently. It is exploratory
and is not on the same footing as Table~\ref{tab6}'s census results; see Sect.~\ref{subsec8_2}.
Generated by \texttt{analyze\_response\_surface\_v2.py}, which prints these bounds so caption and
analysis cannot diverge.}\label{fig6}
\end{figure*}

\subsection{Stability on real replicate data}\label{subsec8_3}

Synthetic ground truth cannot tell us whether the real surfaces resemble any of the four families.
So we also score the criteria on real data by re-evaluation: a verdict is \emph{stable} if a second
evaluation of the same policy pair returns it. We avoid calling these ``independent replicates of
the same evaluation'', because the two sides of a pair differ by more than sampling noise. Some
differ in the policy seed set only; others differ in the simulator port or the inference-stack
build, which Sect.~\ref{subsec5_5} shows is itself a parameter of consequence. What the table
therefore measures is consistency under changes of seed, implementation and protocol together ---
which is the practically relevant notion, and a weaker one than replication in the statistical
sense.

\begin{table*}[t]
\caption{Verdict stability over the eight replicate pairs whose two sides enumerate the same
configuration grid. These eight are \emph{not} eight independent experiments: six are policy pairs
drawn from the same four Octo variants on the eggplant census and therefore share policy data, and
the other two are spoon and carrot, so the eight pairs span three tasks. The underlying records are
the older 48-episode pre-fix sweeps with a per-episode bootstrap, not the current census scope.
\emph{Stable} counts agreement of any kind, including agreement to abstain, so the last two columns
separate out what a reader who wants to trust a declaration needs: how many pairs carried a
declaration on at least one replicate, and how many of those had both replicates declare the same
ordering. At most one declaration successfully replicated under any criterion
(\texttt{analyze\_declaration\_replication.py}).}\label{tab8}
\begin{tabular}{@{}lrrrrrrr@{}}
\toprule
 & & & & & & \multicolumn{2}{c}{Declaring pairs} \\
\cmidrule(l){7-8}
Criterion & Pairs & Stable & Contradict & Unsupported & Declare & $n$ & replicated \\
\midrule
Point calibration & 8 & 0.50 & 0.00 & 0.50 & 0.38 & 5 & 1 (0.20) \\
Union bound, controller & 8 & 0.88 & 0.00 & 0.12 & 0.06 & 1 & 0 (0.00) \\
Union bound, all axes & 8 & 1.00 & 0.00 & 0.00 & \textbf{0.00} & 0 & n/a \\
\bottomrule
\end{tabular}
\end{table*}

The stability column rises from 0.50 to 0.88 to 1.00 down the table, and read alone it would say
that conservatism buys reproducibility. It does, but not in the way that ordering suggests, because
\emph{stable} counts agreement to abstain. The all-axes row reaches 1.00 by \textbf{never declaring
anything} on this data --- the degenerate end of the conservatism dial, where a criterion that always
abstains is perfectly reproducible and perfectly useless. The controller-axis row's 0.88 is seven
abstain--abstain agreements out of eight; it declared on exactly one pair, and that declaration did
\emph{not} reproduce.

This data therefore supports a narrower claim than the column ordering suggests. What we can say is
that the union bound \emph{abstains} more consistently. Point calibration declared on five of the
eight pairs; on four of those five it declared on one replicate and abstained on the other, and the
union bound turns all four into agreed abstentions, which is what drives the unsupported rate from
0.50 down to 0.12. What we cannot say, and do not claim, is that its declarations replicate better. With
one declaration to its name and that one unreplicated, the affirmative rate in the last column
(0 of 1, against point calibration's 1 of 5) is not an estimate of anything; neither number has the
sample size to be ranked against the other. Eight replicate pairs spanning three tasks, six of them
sharing policy data, on a superseded data scope, cannot settle declaration reliability, and we do not
ask them to: our claim about the validity of declarations rests on the synthetic ladder of
Sect.~\ref{subsec8_2}, where ground truth exists and coverage is measured directly. What Table~\ref{tab8} contributes is the trade-off in the middle row --- a real
gain in agreement (0.50 to 0.88) bought with a real loss in declarations (0.38 to 0.06).

\subsection{The delimiting experiment: a pair the real robot reverses}\label{subsec8_4}

Everything to this point is internal to the simulator. The criterion is conservative where we can
check conservatism and stabilises verdicts where we can check stability, but neither test asks
whether it protects against being \emph{wrong about the real robot}. The published tables let that be
asked without hardware, on a case we selected --- and we state the selection, because it is
development information and not a held-out test.

The reference implementation publishes real-robot and simulated success rates for the same policies
on the same task. Auditing all of them
(\texttt{research\_audit/\allowbreak published\_\allowbreak pairwise\_audit.md}), 62
within-task policy pairs have both rates published once RT-2-X is excluded for lack of public
weights, and \textbf{5 of the 62 have opposite mean directions} in simulation and reality, with a
further 3 tied in at least one domain. So this is not a unique case, and an earlier version of this
paper wrongly called it one. We study \texttt{rt-1-converged} vs \texttt{rt-1-15pct} on
pick-coke-can \citep{brohan2023rt1} for three stated reasons: both checkpoints are public, it is on
the embodiment and task where we ran a complete 300-configuration census, and simulation declares it
with a clear margin ($+0.147$). An earlier version of this paper also claimed it has the largest
published real margin of the five. It has the \emph{smallest}: the five real margins are $0.333$,
$0.185$, $0.167$, $0.148$ and, for this pair, $0.067$. That cuts against the case rather than for
it, and the objection it raises is met below rather than omitted. These are reasons of convenience
and coverage, not of significance --- the
five are directional disagreements between \emph{rounded published means}, with no per-cell trial
counts and no intervals on either side, so none of them is an established reversal.

\begin{table*}[t]
\caption{The real-versus-simulation reversal pair. Our reproduction of the published simulated rates
is $+0.006$ and $+0.007$, so the disagreement is reproduced faithfully rather than
introduced.}\label{tab9}
\begin{tabular}{@{}lrrr@{}}
\toprule
Source & rt-1-converged & rt-1-15pct & Margin \\
\midrule
Published real robot & 0.853 & 0.920 & $\mathbf{-0.067}$ (15pct better) \\
Published simulation & 0.857 & 0.710 & $\mathbf{+0.147}$ (converged better) \\
Our nominal census & 0.863 & 0.717 & $+0.147$ \\
\bottomrule
\end{tabular}
\end{table*}

The question is whether the union bound \emph{abstains} on a pair where the simulator's ordering
disagrees with the published real one. It is the sharpest test of the criterion available to us,
because here there is an external number to disagree with --- though that number is a point estimate
without an interval, which we return to below. We ran both checkpoints over the complete
300-configuration census under nominal conditions plus two calibration-invisible conditions --- the
torque limit and object friction --- for 1{,}800 episodes total.

One caveat about that choice of conditions belongs here rather than in a footnote. The replay
evidence of Sect.~\ref{subsec4_3} is on \emph{WidowX} BridgeData demonstrations, and this experiment
is on the Google Robot embodiment. We carried the two conditions over because they are the ones that
moved verdicts on bridge, not because they are established as calibration-invisible for this arm: we
have no Google Robot replay set, so for this embodiment ``calibration-invisible'' is an assumption
imported from the other, and the friction axis is in any case unobserved by free-space replay on
either. The experiment should be read as a test of the criterion's behaviour on a pair with external
ground truth, not as a propagation of a verified compatible set.

The result is adverse. Point calibration on the nominal condition gives $[+0.083, +0.210]$, declaring
\texttt{rt-1-converged} better. The union bound over the two dynamics conditions gives
$L = +0.070$, $U = +0.240$ --- \textbf{also declaring \texttt{rt-1-converged} better}, which is the
ordering the published real-robot rates reverse.

Widening from the point to the set moved the lower bound from $+0.083$ to $+0.070$, so the shift
that would reach abstention is $0.070$: the bound has to lose that much on its lower end, not the
$0.214$ separating the simulated margin from the real one. Measured against $0.070$, the two dynamics
conditions are not close --- their margins span 0.133 to 0.177, a spread of 0.044, which is less than
the distance to zero. But that is a statement about \emph{these two conditions}, not about hidden
variables in general, and the distinction matters: across the twelve texture $\times$ condition cells
the margins span $+0.053$ to $+0.213$, a spread of $0.160$, more than twice what abstention requires.
The axis that carries no physics has the range to move this verdict; the two dynamics parameters do
not.

Having the range to move the verdict is not the same as rescuing the criterion, however, and on
inspection it does not. Because this run puts a second policy on
the same 300-configuration grid, it also settles the question left open by Sect.~\ref{subsec5_5}:
whether a visual variant that carries no physics can flip a \emph{ranking} rather than merely
inflate a single policy's variance. Adding \texttt{urdf\_version} as a fourth set dimension (12 cells
of 75 configurations each) does turn the verdict to abstention, $L = -0.067$ --- but the abstention
is of the second kind identified in Sect.~\ref{subsec7_5}. All twelve cells have \textbf{positive}
point margins, from $+0.053$ to $+0.213$, and not one places \texttt{rt-1-15pct} above
\texttt{rt-1-converged} with interval support. The bound reaches zero only through interval width at
75 configurations per cell.

We put this more carefully than we first did, because the asymptotic statement is easy to overclaim.
Each cell's interval is a policy-noise interval over an enumerated 75-configuration sub-grid, so its
width shrinks with additional \emph{runs} per configuration rather than with additional episodes. If
the twelve observed point margins were the population quantities, $L$ would tend to their minimum,
$+0.053$, and the abstention would dissolve. But they are not: with one run per configuration each
cell margin carries binomial uncertainty we have not resolved, so twelve positive \emph{sample}
means do not establish that all twelve population means are positive, and $+0.053$ is where the
bound tends under our estimand rather than a measured endpoint. What we can say without that
assumption is narrower and enough for the purpose: the abstention here is produced by interval width
and not by any cell disagreeing in sign, so by our own test in Sect.~\ref{subsec7_2} it is at least
partly our power rather than a demonstrated parameter fact, and we decline to claim it as a success
of the criterion. Resolving whether it is \emph{only} power would need repeat runs per configuration
in each cell, which we have not done.

One objection should be met here rather than deferred. The published real margin is $-0.067$,
which is of the same order as the policy-seed noise we measure in simulation, and the published
real rates carry no interval of their own --- so the reversal might be real-robot measurement noise
rather than a model gap. The reading is asymmetric, and either branch supports the same conclusion.
If the reversal is a genuine model gap, the criterion declared an ordering that is wrong. If it is
real-robot noise, then the real evaluation cannot resolve this pair either, and a simulated
declaration still carries no warrant about the robot. What the criterion may not do, on either
branch, is license the sign it produced.

The conclusion we draw is therefore about the criterion's reach, and we bound it to what was
actually tested. The difference between the simulated margin and the published real one is $0.214$,
and one policy's own rate differs by $0.203$ (\texttt{rt-1-15pct}: 0.920 real against 0.717
simulated). Against that, abstention on this pair would have required the union's lower bound to
lose $0.070$. The three dynamics cells we ran span $0.044$ and do not reach it; the twelve
texture $\times$ condition cells span $0.160$ and do reach it, though only by interval width and
without any cell changing sign. So the honest statement has three parts. \emph{Within the settings
we ran}, no configuration of the calibration-invisible parameters produces the real robot's
ordering. \emph{Beyond them} we do not know: we tested two dynamics conditions at one value each and
one visual axis at its four values, out of a parameter family that is continuous in friction,
density, the common gain scale and the torque limit, and we did not explore interactions. And the
comparison target is a point estimate --- the published real rates come without intervals, so the
$0.214$ is a difference between one number with an interval and one without.

We therefore claim only this: on the conditions tested, the criterion declared the ordering that the
published real-robot point estimates reverse. Whether the cause lies outside the whole parameter
family that demonstration replay can constrain is a stronger statement that would need either wider
coverage of that family or a real-robot replication with its own uncertainty, and we do not make it.
In the ledger of Sect.~\ref{subsec3_4} the finding is that $b_{ij}$ is plausibly large on this pair
and that a criterion built on $E_\theta$ and $E_{MC}$ cannot see it --- which is enough to delimit
the criterion, and is all we need it to do.

\subsection{What the ladder is for}\label{subsec8_5}

Taken together the three tests place the criterion precisely, and the placement is narrower than we
set out to establish. It bounds the ambiguity the calibration data leaves in the parameters it can
represent: on synthetic ground truth its observed coverage is conservative on all four surface
families where its smoothness prior holds and fails where it does not
(Sect.~\ref{subsec8_2}); on real replicates it \emph{abstains} consistently where point calibration
declares on one replicate and not the other, which is a claim about abstention and not about
declarations (Sect.~\ref{subsec8_3}); and it changes the verdict on 3 of 17 bridge policy pairs at
the reference protocol's three seed sets, falling to 1 of 17 at five seed sets over a full census
(Sect.~\ref{subsec7_2}). Both counts are uncorrected for multiplicity across pairs; under Holm the
rules disagree on 4 of 17 rather than 1, because the correction removes more of the envelope's
declarations than of point calibration's (Sect.~\ref{subsec7_2}). Which count to quote depends on
whether the claim is about one named pair or about whichever orderings the table declares, and the
extra abstentions are abstentions rather than demonstrated corrections. It does not bound, and cannot bound, the ambiguity arising from a model gap
that no setting of those parameters expresses --- and Sect.~\ref{subsec8_4} exhibits a pair where
reaching abstention would have needed 0.070 while the gap between the simulated and published real
margins is 0.214 and points the other way.

We therefore do not recommend the union bound as a device for deciding which policy is better. We
recommend it as a device for detecting when a benchmark's own evidence cannot decide, which is a
weaker and more defensible role, and for which Sects.~\ref{subsec8_2} and~\ref{subsec8_3} are the
relevant evidence. A declaration from it is not a sim-to-real guarantee; Sect.~\ref{subsec8_4} is the
counterexample that makes that sentence necessary.

%%=============================================================================%%
\section{Limitations}\label{sec9}

We list these in descending order of how much they should change a reader's confidence, rather than
in the order that flatters the paper.

\begin{enumerate}
\item \textbf{The operating-point measurement covers one task and one policy pair.}
Section~\ref{subsec4_4} finds that the replay loss over our grid is minimised at
$(k{\times}2, d{\times}0.5, \text{delay}\ 1)$ on both stacks, and that nominal --- ratio 1, no delay,
where every published rate is computed --- is rejected against that minimiser on our 98
demonstrations. Sect.~\ref{subsec4_3} states why that is a re-calibration result and not a claim
about the original fit: we did not verify that our demonstrations are the trajectories that fit
used, and our search grid may be wider than its.
Section~\ref{subsec7_7} measures what that costs on the eggplant census and finds the point verdict
flipping from abstention to a declaration while the set verdict abstains at both. That is one task,
one policy pair, and one alternative operating point. We do not know whether the pattern --- point
verdicts sensitive to the operating point, set verdicts not --- generalises, and a single instance cannot
establish it; it is the cheapest available test of the question and it is not a survey. Extending it
means the same census at the fitted setting on spoon and carrot, and on more policy pairs.

Within that one task the design is complete rather than partial, and that is worth separating from
the coverage limitation above. Both operating points carry a full six-condition fibre --- twelve
64-configuration censuses in all --- and fibre \emph{membership} at the fitted ratio is measured
rather than assumed: the replay grid varies the common scale at seven other ratios but only once at
$0.25$, so we swept the scale sixteenfold at ratio $0.25$ directly and re-measured the torque limit
at the fitted gains, which does not transfer for free because the fitted point doubles the stiffness
(Sect.~\ref{subsec7_7}). What is narrow here is the number of tasks and pairs, not the design on the
one we ran.

\item \textbf{The criterion's central negative result rests on one selected policy pair.}
Section~\ref{subsec8_4} shows the union bound declaring an ordering the published real means
reverse, on one pair, one task, one embodiment. It is \emph{not} the only such pair: our audit of the
published tables finds 5 of 62 within-task pairs with opposite mean directions, and we chose this one
for public checkpoints, census coverage and a clear simulated margin --- not for the size of its
real margin, which is the smallest of the five. That makes it a case
selected with development information in hand, which is the right way to read it --- a demonstration
that the criterion's reach is limited, not an estimate of how often the limit is reached, and not a
held-out test. The other four are available to anyone wishing to look, and one of them
(\texttt{octo-base} vs \texttt{octo-small} on carrot) is a bridge pair we already have a census for.
A reader should conclude that a union-bound declaration carries no sim-to-real warranty, and should
not conclude that such declarations are usually wrong.

\item \textbf{No hardware of our own.} Every simulated number here is ours; every real-robot number
is quoted from the reference implementation's published tables. Those are the values the benchmark
itself is validated against, and using them is what makes Sect.~\ref{subsec8_4} possible without a
laboratory, but they carry their own uncertainty and are published without intervals: they are
single-protocol real
evaluations, and the pair in Sect.~\ref{subsec8_4} has a real margin of 0.067, of the same order as
the seed noise we measure in simulation.

The path to a real-robot test is short enough to specify rather than gesture at, and we specify it
here as a \emph{plan} whose outcome we do not know. The open question is whether the ordering our
criterion declares on \texttt{rt-1-converged} versus \texttt{rt-1-15pct} holds on the robot; the
published point estimates point the other way, but they carry no interval, so the question is open
in both directions and we do not write it as a known error. The design: those two checkpoints on
pick-coke-can, on hardware matched to the published protocol in the respects that decide an action
--- the same arm, the same control interface and action semantics, the same camera placement and the
same scene layout, since a different 7-DoF arm is not a replication of this benchmark merely by
having seven joints.

Sizing it needs power, not an interval half-width, and the two are easy to confuse. Suppose the
published rates $0.853$ and $0.920$ are the truth, trials are independent Bernoulli, arms are
equally sized, and there is no scene clustering or protocol drift --- all four assumptions
optimistic. Then a two-sided normal-approximation test at $\alpha = 0.05$ has

\begin{table}[h]
\centering
\begin{tabular}{@{}rrrr@{}}
\toprule
per policy & total & expected half-width & power \\
\midrule
200 & 400 & 0.062 & 0.57 \\
400 & 800 & 0.044 & 0.85 \\
800 & 1{,}600 & 0.031 & 0.99 \\
\bottomrule
\end{tabular}
\end{table}

\noindent So $n = 200$ per policy is a coin flip, not a resolution: an expected half-width of $0.062$
below a hypothesised gap of $0.067$ does not make the observed interval exclude zero. About
$400$ per policy --- $800$ trials --- is the point at which the hypothesised gap would be detected
$85\%$ of the time, and an earlier version of this paragraph read the half-width column as though it
were the power column. A paired design, clustering, a different true gap or a different stopping
rule all move these figures, so the table is a magnitude check on the plan and not a formal sample
size (\texttt{scripts/plan\_real\_robot\_trial.py}).

Two things must be frozen before the trials start, because both are choices we have shown decide
answers: the analysis rule (interval, threshold and observation unit) and the simulator operating
point against which the comparison is made. Our release contains the simulated side at census scope,
so the real side is the only missing term. The four other reversal pairs in the 5-of-62 audit of
Sect.~\ref{subsec8_4} extend the same design at the same cost per pair, and one of them ---
\texttt{octo-base} versus \texttt{octo-small} on carrot --- is a bridge pair for which we already
hold a complete configuration census, so it needs only the real arm.

We are stating a protocol we cannot execute: we have no arm, and the two public evaluation services
that would have supplied one were unavailable during this work (Sect.~\ref{subsec1_3}). A reader
with either should be able to run the above from the release without asking us anything.

\item \textbf{One suite, and two kinds of conclusion.} Every measurement here comes from a single
evaluation suite, and the two kinds of claim we draw from it do not travel equally. Claims that
follow from \emph{construction} should hold for any benchmark built the same way: that an
index-driven initial state makes the configuration set a finite population rather than a sample;
that free-space demonstration replay is uninformative about a common gain scaling and about
an inactive saturation limit; that a criterion over a compatible set bounds only the ambiguity that
set can express. Those arguments nowhere depend on which suite, policies or tasks we used. The
\emph{numbers} are specific to this instance and should not be carried anywhere --- 0.121 for the
torque shift, 0.400 for the texture spread, 1 of 17 disagreeing pairs, the 0.038 reproduction. A
reader should take the mechanisms as general and the magnitudes as an existence proof that they can
be large.

\item \textbf{Narrow coverage of policies, embodiments and tasks.} Three model families (two Octo
variants with three deployment configurations, OpenVLA, two RT-1 checkpoints), two embodiments and
four tasks. OpenVLA is run 4-bit quantised, which we did not control against a bf16 baseline; the
torque mechanism in Sect.~\ref{subsec7_4} is a claim about the quantised model as deployed. The
dependence of decision-relevance on the task-policy pair means our coverage is exactly the wrong
shape for extrapolation: we show the effect exists on one task and is nearly absent on another, and
we cannot say which is typical.

\item \textbf{The strong form of the non-identifiability claim is false in our data.} No
calibration-invisible parameter reverses a sign-determined ranking anywhere we looked, and the one
between-condition sign reversal we once claimed did not survive a larger budget
(Sect.~\ref{subsec7_5}).

\item \textbf{The fibre is sampled, not covered.} The union bound evaluates a discrete set of
parameter settings, so its guarantee is of sampling-extremum type, and Sect.~\ref{subsec8_2} shows
it failing (observed coverage 0.22) precisely when the decisive feature lies between sampled
settings. Concretely, we sample the common scale at $\times\tfrac14$ and $\times 4$, the torque
limit at $\times\tfrac12$, and friction and density at one value each --- five points standing in
for a region that is continuous in four coordinates, with no interactions explored. On the contact
axis three sampled points are too few for the GP rung to be trustworthy, so those axes use the union
bound and inherit its limitation.

\item \textbf{Low success rates produce many abstentions, and some are ours rather than the data's.}
Several tasks sit at success rates of 0.1 to 0.5, where even a complete 24-configuration census at
three seeds gives half-widths of 0.05 to 0.13. We applied the $h \to 0$ test wherever the budget
allowed, but for the original-stack rows of Table~\ref{tab6}, which have one run per configuration,
we cannot.

\item \textbf{Implementation-build drift is separated on two policies only.} The $-0.109$ shift and
the 0.73--0.88 episode-level agreement come from a single controlled re-run of two Octo policies
across two inference-stack versions --- one run of each build, so the figure carries no seed
replication and no interval. We treat $\pi$ as an axis of the same kind as $z$ on that basis, which is
thinner evidence than the rest of Sect.~\ref{sec5} rests on.

\item \textbf{Part of the platform work is exploratory.} The ManiSkill3 configuration was used for
the sweeps, with the original reference stack used for cross-checking; the two agree on the replay
invariances to sub-millimetre and on the union-bound verdicts in Table~\ref{tab6}, but they differ on
point verdicts for one task, and their eggplant configuration grids differ so that per-episode
comparison is undefined there. Rows are labelled by stack throughout and should not be pooled.
\end{enumerate}

%%=============================================================================%%
\section{Conclusion}\label{sec10}

This audit connects three properties of a robot-policy evaluation protocol to the comparisons it
supports: the dynamics excited by calibration, the configuration population being evaluated, and
the budget and analysis used to estimate a policy difference. They require different remedies.
Additional independent runs reduce policy randomness, but repeated evaluation at one simulator
setting does not measure sensitivity to other settings or to a different configuration population.

The measurements make this distinction concrete. Replay of 98 demonstrations leaves common gain
scaling and an inactive torque limit weakly constrained over the tested ranges. On a fixed
ManiSkill3 eggplant census, one pair's estimated gap ranges from 0.055 to 0.174 across conditions;
the separate replay bounds in Table~\ref{tab6a} come from the original ManiSkill2 stack. Reapplying
the calibration protocol to same-source demonstrations favours a different controller operating
point on both stacks. A completed census at that point changes one single-setting declaration,
while the paired change in the policy gap remains unresolved and both set envelopes abstain.
This establishes sensitivity under the stated protocol, not an error in the original parameter fit.

The finite configuration grids contain 24 to 300 states, and two ports can enumerate different
populations. A census therefore removes configuration sampling error for its stated population,
not policy randomness or deployment uncertainty. At the published 72-episode budget, one or two
of four reported orderings are resolved, depending on the variance model. Across 17 bridge pairs
at their released budgets, Holm-adjusted single-setting and sampled-set analyses retain nine and
five declarations respectively. The four additional abstentions quantify sensitivity to the sampled
parameter set under the inference assumptions; they are not verified corrections of real-world errors.

These findings motivate a compact reporting protocol. Specify the configuration population, visual
aggregation, simulator and inference builds, observation unit, policy-randomness lifecycle and
replication budget. Report the controller setting used for the rates, its loss on the stated
calibration data, and the best loss found over an explicit search domain. A setting away from that
minimum may be justified by stability, runtime or evidence outside the replay set; reporting these
quantities makes the choice inspectable. Use independent policy replicates to address unresolved
sampling uncertainty, and separate parameter-sensitivity analysis to assess dependence on modelling
choices. A set-valued declaration describes only the finite settings examined. Published real-robot
point estimates motivate further validation, but neither those means nor our simulation audit supply
a guarantee of real-world policy ordering.

%%=============================================================================%%
\bmhead{Acknowledgements}

The authors declare no conflict of interest in connection with this work; the formal statement
appears under Declarations below.

\paragraph{Use of generative artificial intelligence.}
Large language models (Anthropic Claude and OpenAI ChatGPT/Codex) assisted with manuscript drafting
and revision, analysis and figure-generation scripts, and checks of statistical reasoning and
reproducibility. Quantitative results are computed from the recorded episodes by the supplied
scripts. The authors are responsible for the study design, verification and interpretation of the
results and the final manuscript; no language model is listed as an author.

%%=============================================================================%%
\section*{Declarations}

\bmhead{Funding}
%% Confirmed 2026-09-21: no funding of any kind.
The authors declare that no funds, grants, or other support were received during the preparation of
this manuscript.

\bmhead{Competing interests}
The authors have no relevant financial or non-financial interests to disclose.

\bmhead{Ethics approval and consent to participate}
Not applicable. This study involves no human participants and no animals.

\bmhead{Consent for publication}
Not applicable.

\bmhead{Data availability}
Episode-level evaluation records, replay trajectories and analysis inputs are supplied in the frozen
companion review package, with a file manifest, SHA-256 checksums and the scope of the included
experiments. The records identify policy, task, simulator condition, inference build and random
seeds. The package excludes ongoing extensions that do not support a result reported here.
Development is at \url{https://github.com/Jun-Jason-Ji/benchmark-ranking-determinacy}; earlier
archived versions are indexed at \url{https://doi.org/10.5281/zenodo.22893458}. Those earlier
versions are not interchangeable with the current review snapshot, which includes additional
evaluation data. The current snapshot is identified by its accompanying manifest; its
version-specific public archive identifier has not yet been verified. The real-robot reference
values in Sect.~\ref{subsec8_4} come from the benchmark's published source tables, whose location
is cited, rather than from experiments conducted for this study.

\bmhead{Materials availability}
Not applicable.

\bmhead{Code availability}
The companion review package supplies the analysis scripts and reproduction instructions. The
project repository also contains the evaluation harness, configuration-census utilities, resumable
queues and Vulkan compatibility layer described in Appendix~\ref{secB}. Original code is released
under the MIT licence; records and outputs derived from third-party simulators, policies and
demonstrations remain subject to \texttt{NOTICE.md}. The review evidence index maps the main claims
to their records and analysis commands. The current review snapshot should be used when reproducing
the manuscript rather than an earlier archived version.

\bmhead{Author contributions}
%% CRediT roles carried over from the authors' prior submission (IVC_Paper2_Submission_2026-09-12),
%% as directed 2026-09-21. Shengjie Guo is new to this paper and absent from that statement; his
%% roles are those of Xiaolei Zhang, as directed 2026-09-21. Settled -- nothing outstanding here.
\textbf{Jun Ji:} Data curation, Methodology, Formal analysis, Resources, Writing -- original draft.
\textbf{Bowen Tan:} Software, Formal analysis, Writing -- review \& editing.
\textbf{Yi Li:} Supervision, Project administration, Writing -- review \& editing.
\textbf{Xiaolei Zhang:} Supervision, Project administration, Writing -- review \& editing.
\textbf{Shengjie Guo:} Supervision, Project administration, Writing -- review \& editing.
\textbf{Yi Sui:} Supervision, Validation, Writing -- review \& editing.

%%=============================================================================%%
\begin{appendices}

\section{Test inversion for the compatible set}\label{secA}

\subsection{Construction}\label{secA1}

Calibration observes $M$ real demonstrations. For a candidate parameter $z$, let $L_m(z)$ be the
replay error on demonstration $m$ --- the deviation between the simulated and recorded trajectories
--- and let $\hat z$ minimise $\bar L(z) = M^{-1}\sum_m L_m(z)$ over the grid we search.

Rather than report $\hat z$, we invert a one-sided test. For each candidate $z$ we form the paired
per-demonstration differences $\{L_m(z) - L_m(\hat z)\}_{m=1}^{M}$, resample demonstrations with
replacement $B$ times, and take $\mathrm{lo}_{1-\alpha}(z)$ to be the $\alpha$ quantile of the
bootstrap distribution of their mean. Then
\begin{equation}\label{eq:csetA}
C_\alpha = \big\{ z : \mathrm{lo}_{1-\alpha}(z) \le 0 \big\},
\end{equation}
the set of parameters whose replay error is not significantly worse than the best. Pairing by
demonstration matters: the demonstrations differ enormously in difficulty, and an unpaired test
compares that between-demonstration spread against the parameter effect rather than against itself.

The threshold in Eq.~\eqref{eq:csetA} is exactly zero, with no equivalence margin, so \emph{any}
strictly positive lower bound rejects however small it is. That is a consequential choice and we
apply it without exception, including where it costs us: Sect.~\ref{subsec4_4} reports that it
rejects the simulator's own nominal setting on both stacks. \texttt{rebuild\_compatible\_set.py}
exposes the search domain, the loss, $\alpha$ and an optional margin $t \ge 0$ as arguments, and
prints the margin that would be needed to retain any given candidate, so that if a future version of
this analysis wants an engineering-equivalence tolerance it has to name one and defend it rather than
inherit it from a reference choice.

Two properties of this construction carry the argument of Sect.~\ref{sec4}, and we state them with
the measured numbers rather than in the limit.

\paragraph{The insensitive directions move the loss by six to ten orders of magnitude less than the
identified ones.}
One clarification before the numbers, because two different tests appear in this paper and mixing
them is easy. The set construction of Eq.~\eqref{eq:csetA} tests candidates against the grid's
\emph{loss minimiser}, and under that test nominal is rejected (Sect.~\ref{subsec4_4}). The figures
below are a different and narrower comparison: each insensitive direction against \emph{nominal},
which measures the flatness of the loss along that direction and says nothing about whether nominal
itself is retained by the set construction. It is not. What follows is therefore evidence about the
directions, not about the point they pass through.

If $L_m(z) = L_m(z')$ for every $m$, the paired differences vanish identically, the bootstrap
distribution is a point mass at zero, and $\mathrm{lo}_{1-\alpha} = 0 \le 0$ irrespective of $\alpha$,
$B$ or $M$. Our two directions approach that limit rather than attaining it, so what matters is the
measured margin. Tested against nominal on the 98 demonstrations:

\begin{itemize}
\item \emph{Common gain scaling.} Over the fourfold range $\times 0.5$ to $\times 2$ at fixed ratio
and delay, the paired mean loss differences are at most $1.6\times10^{-7}$ with $95\%$ lower bounds
no lower than $-5\times10^{-7}$. Retained.
\item \emph{Torque limit halved.} The paired mean loss difference is $-1.1\times10^{-10}$ on
ManiSkill3 and $-1.8\times10^{-8}$ on the original stack, with exactly one of 98 demonstrations
contributing a nonzero term on each. Retained.
\end{itemize}

\noindent For comparison, the identified directions are rejected with lower bounds of
$+1.5\times10^{-3}$ to $+2.2\times10^{-3}$ (Sect.~\ref{subsec4_4}). Every loss figure in this
appendix is in the composite units of Eq.~\eqref{eq:loss}, which adds metres to radians and is
therefore not a length. The separation between
``retained'' and ``rejected'' here is four to seven orders of magnitude, which is why the conclusion
does not turn on the choice of $\alpha$, of $B$, or on whether one demonstration in 98 behaves
differently. What it does turn on is the reference point, and that is the lesson of
Sect.~\ref{subsec4_4}: tested against the grid minimiser rather than against nominal, the same
procedure retains the same invariant \emph{directions} while rejecting the \emph{point} they pass
through.

\paragraph{Unobserved parameters are compatible for a different reason.}
Contact friction and object density never enter a free-space replay trajectory at all, so their
compatibility reflects the protocol's \emph{coverage} rather than its precision. We flag these
separately throughout, because a reader might otherwise conclude that a longer or better-excited
demonstration set would constrain them. It would not; a different experiment is needed.

\paragraph{Identified directions are excluded with margin.}
Along the gain ratio and the execution delay the same procedure returns strictly positive lower
bounds. Relative to nominal, ratio $\times 0.5$ and $\times 2$ raise the mean paired loss by
$+0.0064$ and $+0.0089$ (Sect.~\ref{subsec4_3}); relative to the grid minimiser, nominal itself is
rejected at $+0.0022$ on ManiSkill3 and $+0.0015$ on the original stack
(Sect.~\ref{subsec4_4}). The construction therefore discriminates, and discriminates sharply enough
to reject the setting the benchmark ships with; it is not vacuously inclusive.

\subsection{The error ledger, and what this paper does not establish}\label{secA2}

Writing $[L_{ij}, U_{ij}]$ for the union-bound interval of Sect.~\ref{subsec3_3}, the claim that
policy $i$ is genuinely better in practice requires three events at once:
\begin{align}
E_\theta &= \{ z^* \in S \}, \\[2pt]
E_{MC}   &= \big\{ \forall z \in S,\ \forall \pi \in \Pi,\ \forall (i,j): \notag \\
         &\qquad \Delta^-_{ij} \le \Delta^{\mathrm{sim}}_{ij}(z; \pi)
                 \le \Delta^+_{ij} \big\}, \\[2pt]
E_b      &= \big\{ \forall (i,j):\ |\Delta^{\mathrm{real}}_{ij} -
             \Delta^{\mathrm{sim}}_{ij}(z^*)| \le b_{ij} \big\},
\end{align}
On $E_\theta \cap E_{MC} \cap E_b$ every real difference lies in
$[L_{ij} - b_{ij},\ U_{ij} + b_{ij}]$ simultaneously, and \emph{if} those events hold at levels
$1-\alpha$, $1-\gamma$ and $1-\beta$, a union bound gives at least $1 - \alpha - \beta - \gamma$
without any independence assumption. We write the levels conditionally because none of the three is
established here --- see the caveats below, of which the first is that our $S$ carries no coverage
guarantee at all. Note that $E_{MC}$ is stated over $\Pi \times S$: the implementation build is
treated as an axis of the same kind as the physical parameters, for the reason given in
Sect.~\ref{subsec5_5}.

Each event needs its own argument, and five caveats are worth stating explicitly because each is a
way the ledger can be misused. The first two describe gaps between what this ledger asks for and what
we report; we do not close either.

\begin{enumerate}
\item \textbf{Neither set we use carries a coverage probability, and for two separate reasons.}
The first is the selection effect of Sect.~\ref{subsec3_2}: $C_\alpha$ as computed picks its
reference from the same demonstrations it then tests on, and under a null of equal-risk candidates
at our $K$ and $M$ that retains a true minimiser about $53\%$ of the time rather than $95\%$. So
even $C_\alpha$ is an empirical rule, not a confidence set; a split-sample variant recovers most of
the level and we report it, but we do not claim exactness for it either. The second is that our
verdicts do not range over $C_\alpha$ at all but over the invisible fibre through nominal
(Sect.~\ref{subsec4_4}), which is generated by a point the calibration data reject. On both counts
$P(E_\theta)$ is unknown for the sets we use, and every ledger statement in this paper is
conditional --- on the benchmark's operating point, and on levels we do not establish.

\item \textbf{Our intervals are per pair; $E_{MC}$ and $E_b$ are joint events.} We compute an
interval for each policy pair separately and report the collection. That is not
$\forall (i,j)$: a family of per-pair statements must be inflated for multiplicity before it can
enter the ledger. Sect.~\ref{subsec7_2} performs that inflation and reports what it costs --- point
calibration 11 declarations to 9 under Holm, the union bound 10 to 5 --- and the knife-edge bounds
of $+0.0091$ and $+0.0095$ are among the ones it removes, as one would expect of margins that small.
What the ledger asks for is still not supplied, because the corrected counts are a family-wise
statement about \emph{these} 17 pairs and the ledger's $E_{MC}$ quantifies over $\Pi \times S$ as
well; but the gap is now measured rather than merely named.
Our per-pair results should be read as exploratory in the multiplicity sense. This is a statement
about \emph{multiple policy pairs}, not about the envelope over conditions: for a fixed finite set
of conditions the probability of a false declaration in either direction is at most
$2\alpha_1 = 0.05$, conditional on valid per-setting tail bounds, without an additional correction
across conditions. The argument in Sect.~\ref{subsec3_3} bounds each pre-specified direction by
$\alpha_1 = 0.025$: a wrong positive declaration requires one particular condition's lower bound
to exceed its truth, and the opposite direction has a symmetric upper-tail event. The two multiplicities are different: across conditions the
error event pins down which interval must fail, across policy pairs it does not.

\item \textbf{$E_\theta$ presumes that $z^*$ exists.} The event is empty of content if no
parameter setting reproduces reality --- that is, if the simulator is structurally misspecified
rather than merely uncalibrated. The fitted parameters may then be no more than locally effective
quantities, and $\alpha$ bounds nothing. Section~\ref{subsec8_4} is best read as a case where this
premise is under strain: no setting of the parameters we can vary reproduces the observed real
ordering.

\item \textbf{A Gaussian-process posterior standard deviation is not $\gamma$.} Substituting the
posterior spread of the response-surface rung into $\gamma$ without validating its calibration
would assume the very smoothness the rung is there to test. Section~\ref{subsec8_2} measures that
rung's coverage on synthetic surfaces instead, including surfaces that violate its prior.

\item \textbf{Surrogate error must be bounded to be absorbed.} Where an approximation has a
certified error bound, the band can be widened by it. Where it does not, that rung is an empirical
diagnostic and must not be reported as a coverage guarantee.
\end{enumerate}

Two further points concern variance reduction rather than coverage. Pairing on common initial states
and common exogenous random draws reduces the variance of $\hat\Delta$ and we use it throughout,
but it does not by itself produce a simultaneous band over $F_{\mathrm{nom}}$ --- simultaneity has to be
constructed, not inherited from pairing. And the internal action sampling of two different policies
does not correspond; forcing it into a pairing and treating the result as a shared exogenous
environment is not valid, which is why Sect.~\ref{subsec5_3} treats the seed as a replication axis
only for policies that consume it.

\textbf{What this paper establishes.} Sections~\ref{sec5}--\ref{sec7} measure $E_{MC}$, and
enumerating $\mathcal{C}$ removes one of its two components exactly. Section~\ref{sec4}
characterises $E_\theta$, whose width dominates the ledger. We do \textbf{not} establish $E_b$: we
have no per-pair sim-to-real bias bound, and Sect.~\ref{subsec8_4} exhibits a pair on which
$b_{ij}$ is large enough to swamp every other term.

\section{Headless rendering of the original stack without a GPU}\label{secB}

The published-protocol reproduction of Sect.~\ref{sec6} and the fractal experiment of
Sect.~\ref{subsec8_4} both require the \emph{original} reference stack rather than the ManiSkill3
port, because the two enumerate different configuration grids (Sect.~\ref{subsec5_1}). That stack
renders through SAPIEN's \texttt{svulkan2}, which requests
\texttt{VK\_KHR\_external\_semaphore\_fd} at device creation. Under WSL2 on Mesa lavapipe no
driver offers file-descriptor-based external semaphores, so device creation fails and the renderer
never initialises --- on a host with no GPU there is no path that satisfies the request.

The extension is requested but, for our workload, unused: fd-based external semaphores serve
cross-API synchronisation, which \texttt{svulkan2} exercises only when CUDA interop is enabled. We
therefore supply an explicit Vulkan layer, \texttt{VK\_LAYER\_the\_world\_fakesemfd}, that makes
the request satisfiable without providing the underlying capability. In 210 lines of C it:

\begin{enumerate}
\item advertises \texttt{VK\_KHR\_external\_semaphore\_fd} and
\texttt{VK\_KHR\_external\_fence\_fd} in \texttt{vkEnumerateDeviceExtensionProperties}, appending
them to whatever the driver reports;
\item strips those two names from \texttt{ppEnabledExtensionNames} at \texttt{vkCreateDevice}, so
the driver is never asked to honour them;
\item removes \texttt{VkExportSemaphoreCreateInfo} and \texttt{VkExportFenceCreateInfo} from
create-info chains; and
\item stubs the fd import and export entry points, which are consequently never reached.
\end{enumerate}

It is an \emph{explicit} layer, enabled per process through \texttt{VK\_LAYER\_PATH} and
\texttt{VK\_INSTANCE\_LAYERS} and installed nowhere:

\begin{verbatim}
export VK_LAYER_PATH=<dir>
export VK_INSTANCE_LAYERS=\
  VK_LAYER_the_world_fakesemfd
\end{verbatim}

\noindent so no other process on the host sees a modified Vulkan environment. It builds with a
single \texttt{gcc -O2 -shared -fPIC} invocation against the Vulkan headers.

\paragraph{Why this does not invalidate the results.}
The layer makes a capability appear present that is absent, which is exactly the kind of
intervention that should be distrusted in a paper about whether evaluation results are determined by
their inputs. Two things bound the risk. First, the mechanism: every path that would consume an
fd-based semaphore is removed rather than faked, so the renderer takes the same code path it would
take on a driver that had never advertised the extension --- the layer does not substitute an
approximate implementation for a real one. Second, and more convincing, the check in
Sect.~\ref{subsec6_2}: this software-rendered, layer-mediated configuration reproduces the
published success rates to a mean absolute difference of 0.038 across eight cells, and to $+0.003$
and $-0.040$ on the second embodiment. A renderer that differed materially from the reference would
not agree to within one seed's worth of noise. The one-signed part of that difference ($-0.038$) is
consistent with a rendering or physics-backend gap and we do not attribute it to the layer
specifically; what the check bounds is the total, which is small.

\paragraph{When it would be unsound.}
If CUDA interop were enabled, the stubbed entry points would be reached and the synchronisation they
are meant to provide would be absent. Any use of this layer with GPU interop active should be
assumed incorrect until re-checked. Our runs use CPU physics with software rendering, so the path is
never taken.

\section{Complete condition table with per-cell counts}\label{secC}

Table~\ref{tabC1} gives, for the pair reported in Fig.~\ref{fig4} and the corresponding rows of
Table~\ref{tab6}, the configurations the task defines, the configurations actually covered, and the
number of observations per configuration. The intervals in the body cannot be read without these: an
interval over a complete census with five runs per configuration and one over a complete census with
two are different objects, and the run count is what distinguishes them. Every cell here is a
complete census, so the configuration sampling error is zero throughout and the intervals carry
policy noise only (Sect.~\ref{subsec5_4}).

The run counts are \emph{not} uniform within a cell, which Remark~\ref{rem:scope}'s shorthand
``$S = 5$'' elides. Where we write $S = 5$ we mean five seed-set directories; the per-configuration
episode count is five on most configurations and four or six on the rest, for the reasons the
caption gives. Any statement of the form ``five runs per configuration'' should be read as the
directory count and checked against this table.

Eq.~\eqref{eq:var} sums over the per-configuration run count, which is the ``Runs/config'' column
rather than the ``Dirs'' column, and Sect.~\ref{subsec7_2} reports the alternative directory unit
and the two rows on which the two differ. The columns here are generated by the same code path as
the tables, so the counts and the estimator cannot disagree.

\begin{table*}[t]
\caption{Per-cell configuration and run counts for octo-small versus octo-base on the three bridge
tasks, under the nominal condition and the five calibration-invisible conditions. ``Configs'' is the
size of the task's configuration grid and ``Covered'' how many were evaluated. ``Dirs'' is the number
of seed-set directories holding records for the cell; ``Min runs'' is the smallest per-configuration
run count in it; and the last two columns give the full distribution the variance estimator uses,
written as runs$\times$configurations. Those three are different numbers, and the reason they differ
is worth giving because it is easy to misread the middle one as the first: seed set B holds
96 episodes on a 64-configuration grid, so ids 64--95 wrap onto configurations 0--31 and give those
an extra observation, while B stops at configuration 47 for three conditions, leaving those
configurations one run short. On the eggplant census the result
is 48 configurations with five runs and 16 with four for three conditions, and 32 with six and 32
with five for two others --- never a uniform five. Eq.~\eqref{eq:var} sums over the per-configuration
$S_c$, so the intervals were computed from the true counts throughout; what was wrong was this
table's summary of them. Generated
from the episode records by \texttt{scripts/make\_appendix\_c.py}.}\label{tabC1}
\centering
\small
% Appendix C, generated by scripts/make_appendix_c.py -- do not hand-edit.
\begin{tabular}{@{}llrrlllll@{}}
\toprule
Task & Condition & Configs & Covered & Dirs & Min runs & Runs/config, S & Runs/config, B & $\Delta$ \\
\midrule
carrot & \texttt{nominal} & 24 & 24 & 2/2 & 2/2 & 2$\times$24 & 2$\times$24 & $-0.042$ \\
 & \texttt{iso\_x0.25} & 24 & 24 & 2/2 & 2/2 & 2$\times$24 & 2$\times$24 & $-0.083$ \\
 & \texttt{iso\_x4.0} & 24 & 24 & 2/2 & 2/2 & 2$\times$24 & 2$\times$24 & $-0.042$ \\
 & \texttt{force\_x0.5} & 24 & 24 & 2/2 & 2/2 & 2$\times$24 & 2$\times$24 & $-0.104$ \\
 & \texttt{fric\_x0.4} & 24 & 24 & 2/2 & 2/2 & 2$\times$24 & 2$\times$24 & $-0.062$ \\
 & \texttt{dens\_x0.5} & 24 & 24 & 2/2 & 2/2 & 2$\times$24 & 2$\times$24 & $-0.042$ \\
\midrule
spoon & \texttt{nominal} & 24 & 24 & 2/2 & 2/2 & 2$\times$24 & 2$\times$24 & $+0.396$ \\
 & \texttt{iso\_x0.25} & 24 & 24 & 2/2 & 2/2 & 2$\times$24 & 2$\times$24 & $+0.250$ \\
 & \texttt{iso\_x4.0} & 24 & 24 & 2/2 & 2/2 & 2$\times$24 & 2$\times$24 & $+0.354$ \\
 & \texttt{force\_x0.5} & 24 & 24 & 2/2 & 2/2 & 2$\times$24 & 2$\times$24 & $+0.375$ \\
 & \texttt{fric\_x0.4} & 24 & 24 & 2/2 & 2/2 & 2$\times$24 & 2$\times$24 & $+0.208$ \\
 & \texttt{dens\_x0.5} & 24 & 24 & 2/2 & 2/2 & 2$\times$24 & 2$\times$24 & $+0.250$ \\
\midrule
eggplant & \texttt{nominal} & 64 & 64 & 5/5 & 5/5 & 6$\times$32, 5$\times$32 & 6$\times$32, 5$\times$32 & $+0.055$ \\
 & \texttt{iso\_x0.25} & 64 & 64 & 5/5 & 4/4 & 5$\times$48, 4$\times$16 & 5$\times$48, 4$\times$16 & $+0.092$ \\
 & \texttt{iso\_x4.0} & 64 & 64 & 5/5 & 4/4 & 5$\times$48, 4$\times$16 & 5$\times$48, 4$\times$16 & $+0.174$ \\
 & \texttt{force\_x0.5} & 64 & 64 & 5/5 & 5/5 & 5$\times$64 & 5$\times$64 & $+0.106$ \\
 & \texttt{fric\_x0.4} & 64 & 64 & 5/5 & 5/5 & 6$\times$32, 5$\times$32 & 6$\times$32, 5$\times$32 & $+0.077$ \\
 & \texttt{dens\_x0.5} & 64 & 64 & 5/5 & 4/4 & 5$\times$48, 4$\times$16 & 5$\times$48, 4$\times$16 & $+0.128$ \\
\bottomrule
\end{tabular}
\end{table*}

\section{Reproduction steps}\label{secD}

The release archived at the DOI in the Data Availability statement carries a
\texttt{REPRODUCIBILITY.md} that is the authoritative and executable version of this appendix; what
follows is its structure.

\textbf{Regenerating the tables and figures} requires no evaluation runs. Every number in this paper
is computed from the episode records rather than transcribed, so the analysis scripts rebuild the
per-pair verdict table (\texttt{make\_core\_table.py}, with \texttt{--unit} selecting the observation
unit of Sect.~\ref{subsec7_2}), the compatible set on the 98-demonstration grids against the loss
minimiser (\texttt{analyze\_compatible\_set\_v2.py}, and
\texttt{rebuild\_compatible\_set.py --split} for the split-sample variant; the older
\texttt{analyze\_compatible\_set.py} is superseded and its defects are recorded in its own
docstring), the operating-point comparison of Sect.~\ref{subsec7_7}
(\texttt{analyze\_fitted\_point.py}, whose \texttt{--selftest} exercises the completeness gate), the
selection-bias check behind Sect.~\ref{subsec3_2} (\texttt{check\_selection\_bias.py}), the reversal
pair of Sect.~\ref{subsec8_4} (\texttt{analyze\_fractal\_reversal.py}), the per-pair torque shift of
Sect.~\ref{subsec7_3} (\texttt{analyze\_torque\_shift\_s5.py}), the build-drift figures of
Sect.~\ref{subsec5_5} (\texttt{analyze\_platform\_drift\_paired.py}), this appendix's table
(\texttt{make\_appendix\_c.py}), and all six figures, which a companion script re-renders at the
journal's geometry and audits against its width and lettering requirements.

\textbf{Re-running the evaluations} is needed only to reproduce the records themselves. Two
simulator stacks are involved and they are not interchangeable: their configuration grids differ, so
per-episode comparison between them is undefined for the eggplant task (Sect.~\ref{subsec5_1}). The
original reference stack is used for Sects.~\ref{sec6} and~\ref{subsec8_4} and runs headless
without a GPU through the layer of Appendix~\ref{secB}; the ManiSkill3 port is used for the
controller sweeps. Requirements are pinned, and the evaluation queues are resumable: each skips
episode ids already present in its target record file, so an interruption costs only the episode in
flight.

\textbf{Four traps} are worth naming for anyone recomputing. Enumerate the configuration grid
rather than sampling it, since episodes past the grid size are exact repeats. Treat the policy seed
as a replication axis only for policies that consume one. Pair any cross-directory comparison on
shared episode ids: the record directories hold different episode counts, and because the episode
index wraps onto the grid, an unpaired comparison silently contrasts different effective scopes ---
computed that way the build effect of Table~\ref{tab3} changes by roughly a factor of two.

And match the \emph{pseudo-random-number lifecycle}, not only the loop structure. A per-episode
re-seed and a single advancing stream can share an outer loop and still measure different things:
the first makes a run one noise realisation replicated across the grid. On the protocol of
Sect.~\ref{sec6} that choice doubled the across-run spread of a cell's rate (0.052 to 0.109), moved
the mean difference from the published rates from $-0.005$ to $-0.038$, and decided the sign of one
comparison --- the retraction in Sect.~\ref{subsec1_2}. Our inference server therefore reports
which lifecycle it applied on every reset, and the sweep asserts it, because the failure mode is
silent: a server without the fix coerces an absent seed to zero and returns plausible numbers under
the wrong protocol.

\section{The retraction record}\label{secE}

Each of these was claimed in an earlier version of this work and each was withdrawn after we ran the
experiment designed to test it. We give the number that settled each one, because a retraction
without its evidence is not checkable.

\begin{enumerate}
\item \emph{``Eggplant nominal: octo-small beats octo-base.''} Declared from one seed set with a
per-episode bootstrap, interval $[+0.08, +0.32]$. Three same-generation seed sets give $+0.141$,
$-0.031$ and $-0.047$. Withdrawn (Sect.~\ref{subsec5_5}).
\item \emph{``The torque effect replicates across two seed sets.''} The episode index fully
determines the configuration, and the policy in question decodes deterministically, so the two
``seed sets'' were repeat runs of identical rollouts --- record-identical in 48 of 48 episodes.
Withdrawn and replaced by a configuration-census protocol (Sect.~\ref{subsec5_3}).
\item \emph{``Budget cannot explain the point-versus-set disagreements.''} Taking the census from
three to five seed sets resolved two of the three disagreements, and dissolved the one
between-condition sign reversal we had claimed. Withdrawn; Sect.~\ref{sec7} states the weaker
conclusion the data supports, and Sect.~\ref{subsec7_5} gives the dissolved reversal in full.
\item \emph{``A faithful re-run of the published protocol reverses a published ordering.''} Our
reproduction placed octo-base above octo-small on carrot where the published table has the reverse.
It was not faithful: the reference seeds its policy once per run and lets one stream advance, while
we re-seeded at every reset, so all 24 configurations in a run shared one noise realisation. Re-run
under the reference lifecycle, carrot agrees in sign with the published table and no cell disagrees.
Withdrawn (Sect.~\ref{subsec6_3}); the budget claim that motivated it survives in the weaker form
that carrot's published $+0.014$ has an interval of $[-0.039, +0.094]$.
\item \emph{Four ranking-flip candidates at 24 and 48 episodes.} Withdrawn before being claimed,
when adding episodes removed them. They are the reason every effect size in this paper is
census-scope.
\end{enumerate}

Four further corrections were to how a result was described rather than to the result: the
compatible set was called a confidence set when it carries no coverage as computed
(Sect.~\ref{subsec3_2}); the test reference was nominal, which makes nominal's own loss increase
zero by construction and so could never reject it (Sect.~\ref{subsec4_4}); the abstention at both
operating points was read as the criterion protecting against the operating-point choice, which an
envelope over a set that excludes the other point cannot do (Sect.~\ref{subsec7_7}); and the case
of Sect.~\ref{subsec8_4} was called the only real-versus-simulated reversal in the published tables
when the audit finds five. The release's \texttt{CHANGELOG.md} carries the per-version record,
including corrections to units, estimator implementation and documentation that never reached a
stated result.

\end{appendices}

\bibliography{references}

\end{document}
